% =====================================================================
%  Branch exchange and multiplicity of the optimal excitation wavelength
%  in nitrogen-vacancy magnetometry under pressure
%
%  Physical Review Applied (REVTeX 4.2, prapplied)
%
%  Compile:  pdflatex main && bibtex main && pdflatex main && pdflatex main
%
%  This is a THEORY manuscript.  No experimental data exist yet.  Its three
%  results are theorems about the structure of the sensitivity optimum; the
%  Ho kernel enters only as a worked example, and every number taken from it
%  is labelled as conditional on that reconstruction.
%
%  PROVENANCE
%    Frozen record ......... docs/theory/theory_freeze_v3_ho_integrated.md
%      body + A1 + A2 + A3 + errata E1-E4 + sanity closure S1
%    Derivations ........... docs/theory/theory_a1_numerical_execution.md
%                            docs/theory/theory_a2_multiplicity.md
%                            docs/theory/theory_a3_branch_exchange.md
%    Positioning ........... docs/audits/novelty_and_exponent_audit.md
%                            docs/audits/positioning_in_high_pressure_sensing.md
%
%  Every number quoted here is produced by code/ and pinned by the test
%  suite.  Do not edit a number here without editing the test that pins it.
%
%  STRUCTURE.  Reorganized 2026-08-28 against the benchmark, from 11 sections
%  + 2 appendices to 8 + 3.  Theorem G now precedes Theorem M because it
%  supplies M's antecedent En > 1; the letters G, M, X are mnemonics, not
%  ordinals.  The old Sec. III dissolved three ways, the phase diagram folded
%  into Sec. V, and the kernel's domain of validity became Appendix B.
%    I    Introduction
%    II   Framework and the two optima ...... (+ old Sec. III A, III B)
%    III  The response exponent ............. Theorem G (A2)
%    IV   Level set and multiplicity ladder .. Theorem M (A2); + the invariance
%         subsection moved out of Sec. III on 2026-08-29
%    V    Pressure-driven branch exchange .... Theorem X (A3, E3) + phase plane
%    VI   Pre-specified tests
%    VII  Discussion ...................... (+ old Sec. III C)
%    VIII Conclusions
%    App. A  Fixed-power optical limit at 120 GPa
%    App. B  Domain of validity of the kernel ...... was Sec. IX (S1, E1, E3)
%    App. C  Reproducing the numbers
%
%  REVIEW OF 2026-08-28.  Findings and the full agenda are in
%  archive/reviews/2026-08-28-main-tex-vs-prapplied-benchmark.md, benchmarked against
%  Sekiguchi et al., Phys. Rev. Applied 21, 064010 (2024).  All four tiers of
%  that agenda are applied: correctness, framing and citations, structure and
%  prose, and numbers and floats.  What was deliberately NOT taken from the
%  benchmark: its 9-page length and its zero tables, both of which follow from
%  its being an experimental report of one number.
%
%  FIGURES.  All five are included as vector PDF with TrueType (not Type 3)
%  fonts, which is what the APS and arXiv submission checkers require, and
%  all five are full-width figure* floats.
%
%  Reworked 2026-08-30, for legibility at print size.  A figure* is included
%  at \textwidth ~= 510 pt = 7.06 in, so a figure drawn at any other width is
%  rescaled, lettering included.  The kernel figure was one 6-panel 15.6 in
%  canvas, i.e. a factor 0.46: its 10 pt labels reached the page at 4.6 pt and
%  its legends at 3.9 pt.  It is now SPLIT IN TWO, the cross-figure check
%  (App. B 1) and the internal checks K1-K4 (App. B 2), each drawn at
%  7.05 in.  The other three were already near 1:1 but carried 5.2-5.9 pt
%  legends and annotations.  code/fig_style.py now holds the one style every
%  figure script uses: 8.5 pt labels, 8 pt ticks, 7 pt legends and in-axes
%  text, and the two canonical widths.  Do not set font sizes in a figure
%  script; change fig_style.py.
%
%  The claim sentences that used to be rendered into the raster now live in
%  the captions, where production can copy-edit them -- this time actually
%  removed from the scripts, which the 2026-08-28 pass had not done.  The
%  per-rung numbers that used to fill the legend of Fig. 2(b) are now the k
%  column of Table III, and the panel carries only the numeral k.
%
%  Regenerate with `python fig_*.py` and `python esa_figv_kernel_sanity.py`
%  in code/ after any change to the figure scripts.
%
%  KNOWN OPEN ITEMS (deliberately not asserted in the text)
%    1. Erratum E4 of the freeze reports that correcting two v1 constants
%       (hbar-omega 65 -> 101.1 meV, dE_ZPL 0.400 -> 0.464 eV) makes v1 pass
%       the Fig. 5(b) gates and puts its optimum near the reconstructed-kernel
%       value.  The
%       script that produces those numbers, code/v1_diagnosis.py, is absent
%       from the repository and from all three archived bundles, so the
%       result is currently unreproducible and is NOT claimed here.  Restore
%       the script, then consider adding the robustness statement to
%       Sec. II D.  NOTE: E4.3's OTHER claim -- one interior maximum for the
%       single-mode envelope against four for the kernel -- was also
%       unreproducible and IS asserted in Sec. II F and Sec. V C.  It is now
%       carried by theory_a1_generalization.v1_envelope_local_maxima and
%       pinned by a test: exactly one maximum at every pressure from ambient
%       to 120 GPa over the full band.
%    2. The zero-phonon critical VALUE a = 0.6938 is a clipped peak height
%       (erratum E3, check K3).  Sec. IV D states the first ladder rung as
%       the interval 1.28 <= I/I_c <= 1.44.  Do not collapse it back to a
%       single number without new source data.
%
%       REVISED 2026-08-30.  This interval used to be 1.44 <= I/I_c <= 2.06,
%       and Sec. IV D used to carry a second, REORDERED ladder for its upper
%       end, in which the zero-phonon step fell past the blue edge and an
%       N = 1 plateau opened.  That branch came from a K3 ratio rising 0.96
%       -> 1.43, which was an artifact of a raster extraction of the source's
%       panel (c) with a constant baseline offset.  The vector extraction
%       (code/extract_ho_fig1_panels_bc.py) puts K3 at 0.88-0.91, FLAT, i.e.
%       the kernel ZPL is ~11 % LOW, not 43 % high.  The correction is an
%       inflation by 1.12, a(ZPL) 0.694 -> 0.780, and 0.780 is still the
%       largest critical value, so the ordering does not change and the
%       sequence 2-4-6-4-3-5-3 holds at BOTH ends of the interval.  The
%       deflated ladder no longer arises and its equation was removed.
%    3. Neither end of that interval is a bound: K3 measures an AREA
%       offset while K2 shows the zero-phonon WIDTH is resolution
%       limited, so the transfer to a height is not justified.  Stated in
%       Sec. IV D.
%    4. Sec. V A's branch pressure coefficients come from one reconstruction
%       of one source.  Lyapin et al. (2018), now cited, is the available
%       external comparison; it has NOT been made.
%    5. refs.bib still carries 12 uncited entries.  Cite or cut before
%       submission.
% =====================================================================
\documentclass[
  aps,
  prapplied,
  reprint,
  amsmath,amssymb,
  floatfix,
  superscriptaddress,
  nofootinbib
]{revtex4-2}

\usepackage{graphicx}
\usepackage{bm}
\usepackage{siunitx}
\usepackage[colorlinks=true,allcolors=blue]{hyperref}

\sisetup{range-phrase=--,range-units=single}

% Float placement.  The earlier version of this manuscript relaxed the budget
% hard (topfraction 0.90, textfraction 0.07) because its three single-column
% figures were 5.5 in tall and could not head a column.  They are now full-width
% and short, so only a mild relaxation is needed.
\setcounter{topnumber}{3}
\setcounter{bottomnumber}{2}
\setcounter{totalnumber}{4}
\renewcommand{\topfraction}{0.85}
\renewcommand{\bottomfraction}{0.50}
\renewcommand{\textfraction}{0.10}

\newcommand{\NVm}{\ensuremath{\mathrm{NV}^{-}}}
\newcommand{\NVz}{\ensuremath{\mathrm{NV}^{0}}}
\newcommand{\gsA}{\ensuremath{^{3}\!A_{2}}}
\newcommand{\esE}{\ensuremath{^{3}\!E}}
\newcommand{\sabs}{\ensuremath{\sigma_{\mathrm{abs}}}}
\newcommand{\Gp}{\ensuremath{\Gamma_{p}}}
\newcommand{\Gps}{\ensuremath{\Gamma_{p}^{*}}}
\newcommand{\lPL}{\ensuremath{\lambda_{\mathrm{PL}}}}
\newcommand{\leta}{\ensuremath{\lambda_{\eta}}}
\newcommand{\labs}{\ensuremath{\lambda_{\mathrm{abs}}}}
\newcommand{\Lset}{\ensuremath{\Lambda_{\eta}}}
\newcommand{\DWF}{\ensuremath{\mathrm{DWF}}}
\newcommand{\Sabs}{\ensuremath{S_{\mathrm{abs}}}}
\newcommand{\SB}{\ensuremath{\mathrm{SB}}}
\newcommand{\ZPL}{\ensuremath{\mathrm{ZPL}}}
\newcommand{\Weff}{\ensuremath{W_{\mathrm{eff}}}}

\begin{document}

\title{Branch exchange and multiplicity of the optimal excitation wavelength
in nitrogen-vacancy magnetometry under pressure}

\author{R.~Abe}
\affiliation{Institute of Science Tokyo, Tokyo, Japan}

\date{\today}

\begin{abstract}
Nitrogen-vacancy centers in the culet of a diamond anvil cell are the reference
local probe of megabar magnetism, and their excitation wavelength is chosen
there, as at ambient pressure, to maximize the photoluminescence. That choice
rests on a condition already known to be violated. Writing the lock-in
sensitivity as $\eta = \Delta\nu/(C\sqrt{R})$ and assuming only that wavelength
acts through the optical pumping rate $\Gp \propto I A(\lambda)$, we prove
three results. A Franck-Condon absorption spectrum carries a zero-phonon
and a sideband branch whose weights scale differently with the Huang-Rhys
factor. Compression raises that factor monotonically, so the two cross at most
once and the optimum \emph{exchanges branch discontinuously}, jumping roughly
\SI{70}{\nano\meter} at a pressure set by the excitation bandwidth---a
dependence no other mechanism in the framework produces. The optimum is a
\emph{level set} $\{\lambda: a(\lambda) = I_c/I\}$ of exactly degenerate
wavelengths, rather than a single wavelength. Its cardinality is a staircase in $\ln I$ whose
steps fall at powers fixed by the \emph{relative} absorption spectrum alone, so
the ladder requires neither absolute calibration nor knowledge of which
mechanism operates. The response enters $\eta$ only through
$E = 2c + s + 2w$, built from the power-law exponents of contrast, rate and
linewidth, and the pump nonlinearity $n$. Splitting is exactly $En>1$. Within
the common-scale power-law response class, responses sharing $(E,n)$ share the
whole $\eta$ surface, so the mechanism cannot be identified from sensitivity
data at any number of powers, and contrast and linewidth must be recorded
separately. The power-dependence model fitted to single nitrogen-vacancy
centers at ambient pressure in the published literature corresponds to
$E = 3$ and $n = 2$. We give eleven pre-specified tests, none
requiring absolute calibration, and state the two spectral regions in which the
reconstructed kernel used for the worked examples must not be used.
\end{abstract}

\maketitle

% =====================================================================
\section{Introduction}
\label{sec:intro}

The search for superconductivity in hydrogen-rich compounds, initiated by
Ashcroft's proposal that metallic hydrogen and hydrogen-dominant alloys should
be high-temperature superconductors~\cite{Ashcroft1968,Ashcroft2004}, has
produced the highest critical temperatures on record: \SI{203}{\kelvin} in
H$_3$S at \SI{155}{\giga\pascal}~\cite{Drozdov2015} and
\SIrange{250}{260}{\kelvin} in LaH$_{10}$ near
\SI{170}{\giga\pascal}~\cite{Drozdov2019,Somayazulu2019}. A second family of
clathrate superhydrides operates at markedly lower pressure: CeH$_9$ is
synthesized at \SIrange{80}{100}{\giga\pascal} and superconducts with $T_c$ up
to \SI{115}{\kelvin} below one megabar~\cite{Salke2019,Chen2021CeH9}, and the
substitutional alloy (La,Ce)H$_9$ reaches \SIrange{148}{178}{\kelvin} over
\SIrange{97}{172}{\giga\pascal}~\cite{Semenok2022LaCeH9}. Establishing that
these materials are superconductors requires local magnetic evidence, and NV
centers implanted a few tens of nanometers beneath the culet supply
it~\cite{Doherty2014,Hsieh2019,Lesik2019,Yip2019}: the sensor is fabricated
inside the anvil, sits micrometers from the sample, and survives the same
pressure. The technique has been pushed from the \SI{30}{\giga\pascal} scale to
\SI{130}{\giga\pascal}~\cite{Hilberer2023}, to megabar
pressures~\cite{Dai2022,Wang2024}, and to
\SI{240}{\giga\pascal}~\cite{Hao2025}, and has imaged the Meissner effect and
flux trapping in CeH$_9$~\cite{Bhattacharyya2024}, in
LaH$_{10}$~\cite{ChenLaH2025}, in conventional superconductors under
pressure~\cite{Dailledouze2025}, and in pressurized
nickelates~\cite{Mandyam2026}.

The scientific content of those experiments is carried by \emph{maps}, and a
map is a raster of optically detected magnetic resonance (ODMR) spectra
whose total acquisition time scales as the
number of pixels times the square of the magnetic
sensitivity~\cite{Barry2020}, at a fixed target field uncertainty. Sensitivity is therefore the currency in which
spatial resolution, field of view, and the number of accessible $(P,T)$ points
are all purchased. One of the choices that sets it---the excitation
wavelength---is normally made by maximizing the detected photoluminescence.
That choice has been optimized before, but never for sensitivity: Beha
\emph{et al.} located an optimum excitation wavelength for single NV$^-$
centers at ambient pressure by balancing excitation against NV$^0$
recharging~\cite{Beha2012}, and Todenhagen and Brandt have measured how the
readout contrast depends on it~\cite{Todenhagen2023}. Both are statements
about the rate or the contrast at one power. Neither is a statement about
$\eta$, neither carries a power axis, and neither has been repeated under
pressure. Reference~\cite{Barry2020} does not treat the excitation wavelength
as an optimization variable at all, taking a single free-running
\SI{532}{\nano\meter} source as sufficient.

That practice rests on a condition, and the condition is testable. The
continuous-wave lock-in sensitivity is $\eta \propto \Delta\nu/(C\sqrt{R})$,
so the wavelength that maximizes the detected rate $R$ minimizes $\eta$ only
if the ratio $C/\Delta\nu$ is stationary there. Todenhagen and Brandt have
measured the excitation-wavelength dependence of the ODMR contrast at ambient
pressure and find it varying between $-12\%$ and $-18\%$ across
\SIrange{480}{575}{\nano\meter}~\cite{Todenhagen2023}. That variation is not
monotone---the same work reports the optical readout performing best between
\SIrange{525}{550}{\nano\meter}---but as a mean logarithmic slope over that
span it corresponds to about \SI{0.43}{\percent\per\nano\meter}. All the
argument below requires is that the slope is nonzero. The condition for the
two optima to coincide is therefore not satisfied, and the question of where
the sensitivity optimum actually sits is open.

Pursuing it does not displace the optimum. It changes what kind of object the
optimum is, and establishing that structure is the purpose of this paper.

The argument rests on a single hypothesis, which we label (M) and treat as an
object of measurement rather than an assumption: that the wavelength enters the
response only through the optical pumping rate $\Gp \propto I\,A(\lambda)$, so
that only the number of absorbed photons matters and not which photon. Under
(M) alone, with no assumption about the shape of the absorption spectrum, three
results follow.

\emph{The response enters through one exponent} (Theorem~G,
Sec.~\ref{sec:G}). For power-law responses the triple $(C,R,\Delta\nu)$ reaches
$\eta$ only through $E = 2c+s+2w$ and the pump nonlinearity $n$. Splitting is
exactly the condition $En>1$. Within that response class, models sharing
$(E,n)$ share the whole $\eta$ surface, so the mechanism is unrecoverable from
sensitivity data at any number of powers.

\emph{The optimum is a level set} (Theorem~M, Sec.~\ref{sec:M}). Given that
interior maximum, the set of optimal wavelengths at power $I$ is
$\{\lambda: a(\lambda) = I_c/I\}$, its members are \emph{exactly} degenerate in
$\eta$, and its cardinality is a staircase in $\ln I$ whose steps sit at ratios
of critical values of the \emph{relative} spectrum---so the whole ladder
follows from a low-power scan, with no absolute calibration and no knowledge of
the mechanism.

\emph{Pressure exchanges the branch} (Theorem~X, Sec.~\ref{sec:X}). A
Franck-Condon absorption spectrum carries a zero-phonon and a sideband branch
whose weights scale differently with the Huang-Rhys factor $S$. Compression
raises $S$, so their ratio crosses at most once, and where it crosses the
global optimum jumps roughly \SI{70}{\nano\meter} between branches, at a
pressure fixed by the excitation bandwidth rather than at a single number.

The two degeneracies have different causes and different signatures---the
ladder is driven by power and needs $I > I_c$, the exchange is driven by
pressure and exists at zero power---so they are separable in practice, and
together they organize the $(P,I)$ plane (Sec.~\ref{sec:phase}). The role of
pressure is thereby
inverted relative to the usual framing of high-pressure sensing: rather than
only a hostile environment a measurement must survive, it is the one
continuously tunable variable that sweeps $\mathrm{d}S/\mathrm{d}P > 0$, and
therefore the only accessible knob that can drive a branch exchange at all.
Temperature changes the linewidths but barely changes $S$.

The worked examples use an absorption kernel reconstructed from the combined
\emph{ab initio} and diamond-anvil-cell study of Ref.~\cite{Ho2026}. The
theorems are statements about the structure of $\eta = 1/(G\sqrt{R})$ and
assume nothing about the spectrum; every number attached to them is conditional
on that reconstruction. Appendix~\ref{sec:kernel} states where it has been
validated and the two spectral regions in which it must not be used.

% =====================================================================
\section{Framework and the two optima}
\label{sec:framework}

This section fixes the two objects the rest of the paper compares---the
wavelength that maximizes the photoluminescence and the wavelength that
minimizes the sensitivity---states the condition under which they
coincide, and states the one hypothesis, (M), on which the theorems rest.
It closes with the absorption kernel used for the worked examples and
with the two spectral regions in which that kernel may not be used.

\subsection{The sensitivity functional}
\label{sec:framework:eta}

The scaling that follows was established for the nitrogen-vacancy
magnetometer by Taylor \emph{et al.}~\cite{Taylor2008}. For continuous-wave,
lock-in-detected ODMR the shot-noise-limited magnetic
sensitivity is~\cite{Barry2020,Dreau2011,Rondin2014}
\begin{equation}
  \eta_B \;=\; \frac{4}{3\sqrt{3}}\,\frac{h}{g_e\mu_B}\,
  \frac{\Delta\nu}{C\sqrt{R}} ,
  \label{eq:etaB}
\end{equation}
with $\Delta\nu$ the ODMR linewidth, $C$ the resonance contrast and $R$ the
detected photon rate. Because we compare wavelengths rather than predict an
absolute noise floor, we work throughout with the reduced quantity
\begin{equation}
  \eta \;\equiv\; \frac{\Delta\nu}{C\sqrt{R}}
  \;=\; \frac{1}{G\sqrt{R}} ,
  \qquad
  G \equiv \frac{C}{\Delta\nu} ,
  \label{eq:eta}
\end{equation}
and with its inverse square, which is the quantity all three theorems are
statements about,
\begin{equation}
  \Phi \;\equiv\; \eta^{-2} \;=\; G^2 R \;=\; \frac{C^2 R}{\Delta\nu^2} .
  \label{eq:Phi}
\end{equation}
Every overall scale---collection efficiency, NV density, integration
bandwidth---cancels from ratios of $\eta$, and none of the results below
depends on the prefactor of Eq.~\eqref{eq:etaB}.

Equation~\eqref{eq:eta} does two things at once. The contrast enters
linearly while the rate enters as a square root, so the two are not
interchangeable and a mechanism acting on $C$ costs four times more than one
acting on $R$ by the same factor. And $\eta$ is a function of the triple
$(C,R,\Delta\nu)$ only through the single combination $\Phi$, which is the
origin of Theorem~G.

\subsection{The fixed-power figure of merit}
\label{sec:framework:A}

At fixed pressure and fixed \emph{incident optical power}, the photon flux is
$I/E_\gamma \propto \lambda I$, so the rate at which photons are absorbed is
proportional to
\begin{equation}
  A(\lambda) \;\equiv\; \lambda\,\sabs\!\left(hc/\lambda,\,P\right) ,
  \qquad
  a(\lambda) \;\equiv\; \frac{A(\lambda)}{A_{\max}} ,
  \label{eq:A}
\end{equation}
and the optical pumping rate is $\Gp = \gamma\,I\,A(\lambda)$ with $\gamma$ a
constant. The distinction between fixed incident power and fixed incident photon flux is
not cosmetic: the two conventions weight the branches of Sec.~\ref{sec:X}
differently by the factor $\lambda_\SB/\lambda_\ZPL$ and therefore return
different exchange pressures. The convention is stated wherever a pressure is
quoted.

We write $\ell_f \equiv \mathrm{d}\ln f/\mathrm{d}\lambda$ for logarithmic
slopes and distinguish three wavelengths,
\begin{align}
  \labs &= \arg\max A , \qquad \lPL = \arg\max R , \notag \\
  \leta &= \arg\min \eta \;=\; \arg\max \Phi .
  \label{eq:three}
\end{align}
Only $\lPL$ and $\leta$ are measurable, and they differ in what must be
recorded: $\lPL$ needs the count rate alone, whereas $\leta$ needs rate,
contrast and linewidth \emph{separately}. Section~\ref{sec:G} shows that this
is a requirement of principle and not of convenience.

\subsection{The coincidence condition}
\label{sec:coincidence:P1}

Since $\eta = 1/(G\sqrt{R})$, the sensitivity optimum maximizes $\Phi = G^2R$
while the photoluminescence optimum maximizes $R$. Write $\ell_X$ for the
logarithmic wavelength derivative of $X$, so that $\ell_\Phi = 2\ell_G+\ell_R$.
Provided $\lPL$ is interior to the analysis window and $R$ and $G$ are
differentiable there, $\ell_R(\lPL) = 0$, and hence
\begin{subequations}
\label{eq:P1}
\begin{equation}
  \text{$\lPL$ is stationary for $\Phi$}
  \quad\Longleftrightarrow\quad
  \ell_G(\lPL) = 0 .
  \label{eq:P1a}
\end{equation}
Stationarity is not optimality, and Eq.~\eqref{eq:P1a} is therefore not a
statement about $\leta$. What ties the two optima together is the stronger
hypothesis that $\ell_G$ vanishes across the window,
\begin{equation}
  \ell_G \equiv 0 \text{ on the window}
  \quad\Longrightarrow\quad
  \leta = \lPL ,
  \label{eq:P1b}
\end{equation}
\end{subequations}
and then they coincide \emph{exactly}, at every pressure and every power, no
matter how strongly the charge yield or the saturation scale depend on
wavelength. The
contrapositive of Eq.~\eqref{eq:P1b} is what an experiment uses: an observed
separation between $\lPL$ and $\leta$ is direct evidence of a
wavelength-dependent $C/\Delta\nu$ and cannot be produced by the detected rate
alone.

The distinction between stationarity and optimality is not academic. Case~(c)
of Sec.~\ref{sec:coincidence:P3} is a counterexample to the converse of
Eq.~\eqref{eq:P1a}: under (M) the photoluminescence optimum sits at an interior
maximum of $A$, so $\ell_G(\lPL) = 0$ holds there automatically, and yet above
the threshold power $\lPL$ is a local \emph{minimum} of $\Phi$ and is not a
sensitivity-optimal wavelength at all. Nor need
the separation be small where it occurs. On the reconstructed kernel of
Sec.~\ref{sec:framework:kernel} at \SI{120}{\giga\pascal}, a
wavelength-independent slope $\ell_G$ above
\SI{0.27}{\percent\per\nano\meter} carries the global optimum from the sideband
maximum at approximately \SI{441}{\nano\meter} to the zero-phonon line at
\SI{514.5}{\nano\meter}, a displacement of \SI{73.8}{\nano\meter}; that
threshold is conditional on the reconstruction, but its order of magnitude is
not.

Equation~\eqref{eq:P1b} is what makes the measurement of
Ref.~\cite{Todenhagen2023} load bearing: the contrast slope quoted in
Sec.~\ref{sec:intro}, $\ell_C \approx \SI{0.43}{\percent\per\nano\meter}$,
gives $\ell_G \neq 0$ unless the linewidth happens to carry the same slope, and
the two optima are then distinct. That proviso is not idle---no
wavelength-resolved linewidth has been published---which is why $\ell_G$ is the
object of test~T3 rather than an input to the framework.

\subsection{The mediation hypothesis (M)}
\label{sec:framework:M}

All three theorems are conditional on one hypothesis.

\begin{quote}
\textbf{(M)} The response functions---the \NVm{} fraction and detection
efficiency $q_-$, the contrast $C$, the linewidth $\Delta\nu$, and the
saturation scale---depend on the excitation wavelength \emph{only} through the
optical pumping rate $\Gp \propto I\,A(\lambda)$; that is, only on how many
photons are absorbed and not on which photon.
\end{quote}

Under (M) the sensitivity collapses to a function of a single variable,
$\eta = h(\Gp)$, and the wavelength enters only through $A$. This is what makes
the level-set structure of Theorem~M possible, and it is why we treat (M) as
the object of the experimental program rather than as a background
assumption.

(M) does not say that the response is trivial. It is compatible with
arbitrarily strong wavelength dependence of $q_-$, $C$ and $\Delta\nu$, so long
as that dependence is mediated by $\Gp$. It fails only for a channel that
depends on the photon \emph{energy} directly---an ionization edge, a resonance
in the conduction band, an absorption band belonging to a different charge
state whose envelope differs in shape from that of \NVm{}.

(M) is known to fail somewhere. Reference~\cite{Todenhagen2023} reports sharp
resonant features in the electrically detected magnetic resonance at
\SI{521}{\nano\meter}, assigned to a resonant level in the conduction band, and
at \SI{575}{\nano\meter}, the \NVz{} zero-phonon line, where it records the
largest contrast magnitude it observes. Both are photon-energy specific and
neither can be produced by absorbed photon number alone. The charge-state
channel that carries them also has a continuous wavelength dependence, measured
in Ref.~\cite{Aslam2013} and calculated from first principles in
Ref.~\cite{Razinkovas2021}, so (M) is an approximation everywhere
and not only at the resonances. Both features, however, lie redward of
\SI{500}{\nano\meter}, outside the blue window that carries the high-pressure
optimum, and the same reference states that the contrast is otherwise rate
limited by the same process as the photoluminescence excitation. On the present
evidence the data are consistent with (M) away from the identified resonances
and inconsistent with it at them---exactly the situation in which a null test
of arbitrary precision is valuable. Test~T6 of Sec.~\ref{sec:tests} is that test.

Finally, the assumption that $\Gp$ is \emph{linear} in $I A(\lambda)$ can be
dropped. If $\Gp$ is any monotone function of $\xi = I A(\lambda)$---as it is
in Ref.~\cite{Dreau2011}, where the pumping rate saturates---then
$\eta = h(\Gp(\xi)) = \tilde h(\xi)$, the optimal set is still
$\{\lambda: A(\lambda) = \xi^*/I\}$, and the transition-power ratios of
Theorem~M are unchanged. Only the conversion from $\xi$ to laser power becomes
nonlinear.

\subsection{Three cases under the mediation hypothesis}
\label{sec:coincidence:P3}

Under (M), $\eta = h(\Gp)$ and $\Gp$ depends on $\lambda$ only through $A$.
Three cases exhaust them.

\begin{enumerate}
  \item[(a)] If $R$ is increasing in $\Gp$ throughout, then
  $\lPL = \labs$ at all powers.
  \item[(b)] If $\Phi$ is increasing in $\Gp$ throughout, then
  $\leta = \lPL = \labs$: the conventional practice is correct and there is
  nothing to report.
  \item[(c)] If $\Phi$ has an interior maximum at
  $\Gps < \gamma I A_{\max}$, then $\leta$ is \emph{not a point}. It is the
  solution set of $A(\lambda) = \Gps/(\gamma I)$, whose members satisfy
  $\eta(\lambda_i) = \eta(\lambda_j)$ exactly. Note that $\labs$ is
  \emph{not} one of them: the level $\Gps/(\gamma I)$ lies strictly below
  $A_{\max}$ in this case, so the solution set brackets $\labs$ without
  containing it. At $\labs$ the pumping rate is maximal and therefore past
  $\Gps$, where $\Phi$ is decreasing in $\Gp$; $\labs$ is thus a stationary
  point of $\Phi$ but a local \emph{minimum} of it, and the sensitivity there
  is strictly worse than at any member of the optimal set.
\end{enumerate}

Case (c) is the subject of Theorem~M. The degeneracy in (c) is exact---we
verify it numerically to $2.2\times10^{-16}$---but it is not symmetric on the
wavelength axis. On the reconstructed kernel the red member sits about
\SI{3}{\nano\meter} further from $\labs$ than the blue member at every level,
because the two flanks carry different curvature. Correspondingly the $5\%$
penalty band at \SI{120}{\giga\pascal} is \SIrange{426.4}{457.9}{\nano\meter}
in the optical limit of Appendix~\ref{app:optical},
which is approximately $441^{+17}_{-15}$, an asymmetric interval. A test that
searches for members placed symmetrically about $\lPL$ will miss them; the test
must search for equal $\eta$.

\subsection{The optical kernel is taken from outside}
\label{sec:framework:kernel}

The theorems require no model of $\sabs$. The worked examples do, and we take
it from Ref.~\cite{Ho2026}, whose combined \emph{ab initio} and DAC study
reports pressure-dependent absorption spectra, Huang-Rhys~\cite{Huang1950}
factors and Debye-Waller factors to \SI{120}{\giga\pascal}, by the vibronic
method of Ref.~\cite{Razinkovas2021vib}. At ambient pressure those quantities
are to be compared with the first-principles lineshape of
Ref.~\cite{Alkauskas2014}, which treats the emission side, and with the
measured vibronic band at \SI{1.945}{\electronvolt} of
Ref.~\cite{DaviesHamer1976}; we make no quantitative comparison here. The
kernel is a reconstruction of the published
absorption curves; it is not an independent calculation, and we make no claim
to one.

One fact forces that choice, and it is the fact on which the numerical
content of Theorem~M rests. A single-effective-mode Franck-Condon
envelope has exactly \emph{one} interior local maximum, at every pressure from
ambient to \SI{120}{\giga\pascal}, and never splits. The
reconstructed kernel has \emph{four} at \SI{120}{\giga\pascal}
(Table~\ref{tab:maxima}). Every rung of the multiplicity ladder lives on the
critical points and edges of such a kernel, so no single-mode envelope can
produce them at any parameter value. This is a statement about structure, not about goodness of fit, and it
should not be argued from a fit score. Appendix~\ref{sec:kernel} sets out where
the reconstruction has been validated and where it has not.

\begin{table}[tb]
  \caption{The fixed-power figure of merit $A(\lambda) = \lambda\sabs$ of the
  reconstructed kernel at \SI{120}{\giga\pascal}, normalized to its maximum,
  and the optical-limit sensitivity penalty $a^{-1/2}$ each wavelength would
  incur. The upper block lists the interior \emph{maxima} of $A$; rungs of the
  multiplicity ladder sit on these, on the two interior minima and on the
  window edge, for which Table~\ref{tab:ladder} gives the full set. The lower
  block lists commercial laser lines. These values are conditional on the
  reconstruction; for Theorem~M what matters is the \emph{existence} of several
  maxima rather than the values themselves.
  Positions carry a systematic $\pm\SI{1}{\nano\meter}$: an independent trace
  of the same source figure places the sideband maximum
  \SI{0.85}{\nano\meter} away. Values of $a$ carry the $1\%$ height agreement
  of Appendix~\ref{sec:kernel:cross}. The zero-phonon row is the only one
  resting on a peak \emph{height} rather than a position, and carries in
  addition the constant $11\%$ offset against the published Debye-Waller
  factor measured by check~K3 of Appendix~\ref{sec:kernel:internal}. The \SI{532}{\nano\meter} row is an
  interpolation artifact quoted only to be excluded
  (Appendix~\ref{sec:kernel:forbidden}).}
  \label{tab:maxima}
  \begin{ruledtabular}
  \begin{tabular}{lccl}
    $\lambda$ (nm) & $a = A/A_{\max}$ & penalty $a^{-1/2}$ & assignment \\
    \colrule
    \multicolumn{4}{l}{\emph{interior maxima of $A$}} \\
    440.6\textsuperscript{a} & 1 & 1 & sideband, global max \\
    475.6 & 0.660 & 1.231 & sideband \\
    500.2 & 0.277 & 1.899 & sideband \\
    514.5 & 0.694--0.780 & 1.13--1.20 & zero-phonon line \\
    \colrule
    \multicolumn{4}{l}{\emph{commercial laser lines}} \\
    405 & 0.586 & 1.306 & at the ionization edge \\
    445 & 0.992 & 1.004 & inside the $5\%$ band \\
    457 & 0.916 & 1.045 & inside the $5\%$ band \\
    473 & 0.688 & 1.205 & outside \\
    488 & 0.423 & 1.538 & outside \\
    505 & 0.110 & 3.02 & outside \\
    532 & 0.006 & 13 & interpolated only \\
  \end{tabular}
  \end{ruledtabular}
  \begin{flushleft}\footnotesize
  \textsuperscript{a}Numerical maximum of the interpolant; reported in prose as
  approximately \SI{441}{\nano\meter}, not as sub-nanometre source precision.
  \end{flushleft}
\end{table}

% =====================================================================
\section{The response exponent}
\label{sec:G}

Theorem~M, proved in the next section, needs an antecedent: the
sensitivity must have an interior optimum in the pumping rate at all.
This section supplies it, and finds that the antecedent is both easier to
meet and more consequential than it looks. Easier, because the whole response
enters through one number, so a single measured exponent decides it. More
consequential, because that same collapse makes the mechanism unrecoverable
from the sensitivity.

\begin{quote}
\textbf{Theorem G.} $\eta$ depends on the response triple $(C,R,\Delta\nu)$
only through $\Phi = C^2R/\Delta\nu^2$. For power-law responses
\begin{equation}
  R = \Gp(1+\rho)^{-s}, \quad
  C = (1+\rho)^{-c}, \quad
  \Delta\nu = (1+\rho)^{w} ,
  \label{eq:powerlaw}
\end{equation}
written in terms of the reduced variable
\begin{equation}
  \rho = \left(\Gp/\Gamma\right)^{n} ,
  \label{eq:rho}
\end{equation}
with $\Gamma$ a common response scale and $n$ the exponent with which the
response variable follows the pumping rate, one has
\begin{equation}
  \Phi = \Gp\,(1+\rho)^{-E},
  \qquad
  \boxed{\;E \equiv 2c + s + 2w\;} ,
  \label{eq:E}
\end{equation}
so that $\Phi$ has an interior maximum if and only if $En > 1$, and then
\begin{equation}
  \rho^* = \left(\frac{\Gps}{\Gamma}\right)^{n} = \frac{1}{En-1} .
  \label{eq:rhostar}
\end{equation}
\end{quote}

The proof is a substitution:
$\Phi = (1+\rho)^{-2c}\,\Gp(1+\rho)^{-s}\,(1+\rho)^{-2w} = \Gp(1+\rho)^{-E}$,
whence $\mathrm{d}\ln\Phi/\mathrm{d}\ln\Gp = 1 - En\rho/(1+\rho)$, which
vanishes at $\rho^* = 1/(En-1)$ and only for $En>1$. $\blacksquare$

The response therefore enters $\eta$ through the pair $(E,n)$ and nothing else.
The exponent $n$ is a property of how the response variable is driven rather
than of which response is affected: $n = 1$ is the linear case, and
Sec.~\ref{sec:G:measured} exhibits a published fit with $n = 2$.

Table~\ref{tab:gauge} evaluates $E$ for the standard CW-ODMR mechanisms at
$n = 1$. The content of the table is that the question ``which mechanisms can
split the optima?'' collapses to the sign of $En-1$. Rate saturation alone and
linewidth power broadening alone both give $E = 1$ and cannot split at
$n = 1$. Saturation together with broadening gives $E = 2$ and can. Contrast
collapse alone also gives $E = 2$ and can---because $\eta$ is first order in
$C$ and only half order in $R$, which is the structural asymmetry noted in
Sec.~\ref{sec:framework:eta}. All three together give $E = 4$ and
$\rho^* = 1/3$.

Two further restrictions bound the degeneracy. It holds at fixed $n$: two
mechanisms with the same $E$ but different $n$ have different $\rho^*$ and
different $\Phi$, and at $n = 2$ the rate-saturation row of
Table~\ref{tab:gauge} does split. And it holds within one functional class,
since $\Phi$ sees the whole response and not only its exponents; a response
outside the form of Eq.~\eqref{eq:powerlaw} with the same nominal $(E,n)$ is
not degenerate with these.

The composite rows of Table~\ref{tab:gauge} require the mechanisms to share a common half-scale $\Gamma$;
with distinct scales $E$ is not a single scalar, and the gauge degeneracy of
Sec.~\ref{sec:G:identifiability} is exact only in the regime where one scale
dominates or the scales coincide.

\begin{table}[tb]
  \caption{The splitting exponent $E = 2c + s + 2w$ for the standard CW-ODMR
  response mechanisms, evaluated at $n = 1$ and with a common response scale.
  Splitting---and therefore the entire level-set structure of Theorem~M---is
  exactly the condition $En>1$. The last two columns are what $\eta$ can see;
  the mechanism labels in the first column are what it cannot
  (Sec.~\ref{sec:G:identifiability}). The final row is the published fit
  of Ref.~\cite{Dreau2011}, which carries $n = 2$
  (Sec.~\ref{sec:G:measured}).}
  \label{tab:gauge}
  \begin{ruledtabular}
  \begin{tabular}{lccccccc}
    mechanism & $c$ & $s$ & $w$ & $E$ & $n$ & splits & $\rho^*$ \\
    \colrule
    rate saturation alone     & 0 & 1 & 0          & 1 & 1 & no  & --- \\
    power broadening alone    & 0 & 0 & $1/2$ & 1 & 1 & no  & --- \\
    saturation + broadening   & 0 & 1 & $1/2$ & 2 & 1 & yes & 1 \\
    contrast collapse alone   & 1 & 0 & 0          & 2 & 1 & yes & 1 \\
    all three                 & 1 & 1 & $1/2$ & 4 & 1 & yes & $1/3$ \\
    \colrule
    Ref.~\cite{Dreau2011}     & 1 & 0 & $1/2$ & 3 & 2 & yes & $1/5$ \\
  \end{tabular}
  \end{ruledtabular}
\end{table}

\subsection{The antecedent \texorpdfstring{$En>1$}{En>1} is a measured
quantity}
\label{sec:G:measured}

Theorem~M requires an interior maximum of $\Phi$, and Theorem~G reduces that
requirement to $En>1$. That requirement is not assumed here.
Reference~\cite{Dreau2011} measured $R$, $C$ and $\Delta\nu$ against power for
single NV centers at ambient pressure and fitted them to closed forms; $E$ and
$n$ are read off those forms, and give $En = 6$.

Reference~\cite{Dreau2011} reports, for single NV centers, the closed forms
\begin{equation}
  R = R_\infty x, \quad
  C = \Theta\,\frac{1}{1+\rho}, \quad
  \Delta\nu = \frac{\Gamma_c^\infty}{2\pi}\sqrt{\kappa}\,\sqrt{1+\rho} ,
  \label{eq:dreau}
\end{equation}
in which $\varsigma = P_{\mathrm{opt}}/P_{\mathrm{sat}}$ is the optical pumping
power normalized to saturation, written $\varsigma$ to keep it distinct from
the rate exponent $s$, and $x = \varsigma/(1+\varsigma)$ is the resulting
saturated pumping rate. The remaining symbols are the saturated contrast
$\Theta$ and rate $R_\infty$, the saturated pumping and coherence rates
$\Gamma_p^\infty$ and $\Gamma_c^\infty$, the Rabi frequency $\Omega_R$, and the
combinations $\kappa = \Omega_R^2/(\Gamma_p^\infty\Gamma_c^\infty)$ and
$\rho = x^2/\kappa$.

The identification with Eq.~\eqref{eq:powerlaw} is through $x$, not through
$\varsigma$: it is $x$ that plays the role of $\Gp$, because $R$ is exactly
linear in it. This is the nonlinear-pumping case admitted in
Sec.~\ref{sec:framework:M}, and it is why the exponent $n$ of
Eq.~\eqref{eq:rho} is needed. Reading off the exponents gives $c = 1$,
$s = 0$ and $w = \tfrac12$, so $E = 3$; and because $\rho$ is quadratic in $x$
rather than linear, $n = 2$ and $En = 6$, whence
\begin{equation}
  \rho^* = \frac{1}{En - 1} = \frac{1}{5} .
  \label{eq:rho02}
\end{equation}
The numerical maximizer reproduces that identity to five decimals, at
$x^* = \sqrt{\kappa/5} = 0.067$ and $\varsigma^* = 0.072$, i.e. at $7\%$ of the
saturation power for the representative drive
$\Omega_R = \SI{3e6}{\radian\per\second}$.

The \emph{splitting condition} is drive independent: $\rho^* = 1/5$ is
identical for all five
microwave drive strengths Ref.~\cite{Dreau2011} reports
($\Omega_R = 1.1$, $2.5$, $3.0$, $3.6$ and
$5.5\times10^{6}\,\si{\radian\per\second}$), so it is a property of the
functional form and cannot be removed by microwave tuning. The \emph{power at
which the optimum falls} is not: since $x^* = \sqrt{\kappa/5}$ and
$\kappa \propto \Omega_R^2$, $\varsigma^*$ scales with $\Omega_R$, running from
$0.025$ to $0.14$ across those same five settings. Any statement about where in
the power sweep the optimum sits must carry the drive strength with it.

Reference~\cite{Dreau2011} plots that optimum and describes it in words, so the
interior maximum of $\Phi$ was observed in 2011. What Theorem~M adds is that
under (M) this single fact makes the optimal set a level set, degenerate and
multi-valued, with a ladder fixed by the relative spectrum.

\subsection{The mechanism is not identifiable from the sensitivity}
\label{sec:G:identifiability}

\begin{figure*}[t]
  \includegraphics[width=\textwidth]{figures/fig_g_gauge_degeneracy.pdf}
  \caption{The gauge degeneracy of Theorem~G. (a) The reduced figure of
  merit $\Phi = C^2R/\Delta\nu^2$, normalized to its maximum, against the optical
  pumping rate $\Gp$ in units of the half-scale $\Gamma$, on a logarithmic
  abscissa, for five representatives of the plane $E = 2c+s+2w = 2$ at
  $n = 1$, two of which appear in Table~\ref{tab:gauge}. The five curves are drawn and are
  indistinguishable: the largest relative difference anywhere on the plane is
  $6.5\times10^{-16}$, so no sensitivity measurement at any number of powers
  can separate them. (b) The three individually measurable factors for the
  same five models at $\Gp = \Gamma$, in arbitrary units in which the first
  model has $R = \Delta\nu = 1$: detected rate $R$, contrast $C$ and linewidth
  $\Delta\nu$, which span a factor $2$ across the five. Separating the
  mechanisms requires recording $C$ and $\Delta\nu$ apart; measuring $\eta$
  better cannot do it.}
  \label{fig:gauge}
\end{figure*}

At fixed $n$, the models with a given $E$ form a plane $2c + s + 2w = E$ in
exponent space, and every model on that plane has an identical $\Phi$. We
verify the degeneracy numerically across the plane $E = 2$ at $n = 1$: the
maximum relative difference in $\Phi$ between members is $6.5\times10^{-16}$
[Fig.~\ref{fig:gauge}(a)]. It follows that the entire $\eta(\lambda,I)$
surface, the value of $\Gps$, the critical intensity
$I_c \equiv \Gps/(\gamma A_{\max})$ defined in Eq.~\eqref{eq:levelset} below, and the
whole multiplicity ladder are identical for all of them.

\emph{Mechanism attribution from wavelength scans is therefore impossible, at
any number of powers, within the common-scale power-law response class.} An
observed split does not identify contrast collapse as
its cause, because saturation combined with power broadening produces exactly
the same $\eta$ surface. What separates them is the individual responses at a
given pumping rate: at $\Gp = 1$, contrast collapse gives $(R,C,\Delta\nu) =
(1.000,\,0.500,\,1.000)$ while saturation with broadening gives
$(0.500,\,1.000,\,1.414)$ [Fig.~\ref{fig:gauge}(b)]. This yields a
hierarchy of identifiability:
\begin{center}
\begin{tabular}{ll}
  \hline\hline
  measured & determined \\
  \hline
  $\eta$ alone     & $\Phi = C^2R/\Delta\nu^2$ only (degenerate) \\
  $+\ R$           & $G = C/\Delta\nu$ \\
  $+\ C$ or $\Delta\nu$ & all of $(C,R,\Delta\nu)$ \\
  \hline\hline
\end{tabular}
\end{center}

Recording $C$ and $\Delta\nu$ separately is therefore gauge fixing, not
laboratory hygiene: an experiment that records only the sensitivity, or only
the sensitivity and the rate, \emph{cannot in principle} identify the
mechanism at any precision. None of this reaches the ladder, however: its
transition powers depend on the response only through $\Gps$, so the degeneracy
that defeats mechanism attribution leaves the prediction of the next section
untouched.

% =====================================================================
\section{Level set and multiplicity ladder}
\label{sec:M}

Case~(c) of Sec.~\ref{sec:coincidence:P3} is the interesting one, and
Sec.~\ref{sec:G:measured} has just shown that published data put us in
it. What has never been asked is what the solution set of that case looks
like. It is not a point, and its size is an observable.

\begin{quote}
\textbf{Theorem M.} Assume (M), and let $\Phi$ be continuous with
$\Phi(\Gp) \to 0$ as $\Gp \to 0^+$ and with a \emph{unique interior local
maximum}, at $\Gps$---which by Theorem~G exists exactly when $En > 1$, and
which Sec.~\ref{sec:G:measured} shows is measured. Assume further that $A$ is
continuous and piecewise $C^1$ on the closed analysis window
$[\lambda_b,\lambda_r]$, with finitely many strict local extrema, with its
interior critical values pairwise distinct and distinct from $A(\lambda_b)$ and
$A(\lambda_r)$, and strictly monotone in a neighborhood of each endpoint. Then
the set of sensitivity-optimal wavelengths at incident power $I$
is the level set
\begin{equation}
  \Lset(I) = \left\{\lambda \;:\; a(\lambda) = \frac{I_c}{I}\right\},
  \qquad
  I_c \equiv \frac{\Gps}{\gamma A_{\max}} .
  \label{eq:levelset}
\end{equation}
Its cardinality $N(I) = |\Lset(I)|$ is piecewise constant in $\ln I$ and
changes only where the level $I_c/I$ crosses a critical value of $a$, by
\begin{equation}
  \Delta N =
  \begin{cases}
    +2 , & \text{interior maximum},\\
    -2 , & \text{interior minimum},\\
    -1 , & \text{edge, $A$ increasing inward},\\
    +1 , & \text{edge, $A$ decreasing inward},
  \end{cases}
  \label{eq:deltaN}
\end{equation}
and the transition powers are
\begin{equation}
  \boxed{\;\frac{I_k}{I_c} = \frac{1}{a_k} = \frac{A_{\max}}{A_k}\;}
  \label{eq:rungs}
\end{equation}
with $A_k$ the critical values of $A$.
\end{quote}

\emph{Proof.}
Under (M), $\eta$ is a function $h(\Gp)$ of the pumping rate alone, and $\Gp$
depends on $\lambda$ only through $A(\lambda)$. Hence
\begin{equation}
  \arg\min_\lambda \eta
  = \{\lambda : \Gp(\lambda) = \Gps\}
  = \{\lambda : A(\lambda) = \Gps/(\gamma I)\} ,
\end{equation}
which is Eq.~\eqref{eq:levelset}. A continuous $\Phi$ that vanishes at the
origin and has a single interior local maximum is strictly increasing on
$(0,\Gps)$; hence for $I < I_c$ the required level exceeds $A_{\max}$, is
unattainable, and the constrained minimum is the single point $\labs$. For
$I > I_c$ the level $a^* = I_c/I$ decreases monotonically as $I$ rises.

The counting is elementary and does not need $A$ to be smooth, which matters
because the kernel used later is a piecewise-linear interpolant and so has no
nondegenerate critical points in the sense of Morse theory. Between consecutive
critical values, $A$ is strictly monotone on each of the finitely many
subintervals into which its extrema divide the window, so $A^{-1}(y)$ contains
exactly one point per subinterval that spans $y$, and that count cannot change.
As $y$ descends through a strict interior maximum, the two subintervals meeting
there begin to span $y$ together and two preimages appear; through a strict
interior minimum, two merge and disappear. Pairwise distinct critical values
are what keeps these events separate, so that the observable step is $\pm2$
rather than a superposition. At an endpoint $\lambda_e$, where $A$ is monotone,
a single branch exists for $y$ on one side of $A(\lambda_e)$ and leaves the
domain on the other: one preimage is lost if $A$ increases into the interior
and one is gained if it decreases, which is the last two lines of
Eq.~\eqref{eq:deltaN}. Both edges of the window used in
Sec.~\ref{sec:M:example} increase inward, so both contribute $-1$ there.
Setting $a^* = a_k$ gives $I_k = I_c/a_k$, which is
Eq.~\eqref{eq:rungs}. $\blacksquare$

For a smooth $A$ the same count is the regular interval theorem together with
the Morse lemma~\cite[Thm.~3.1]{Milnor1963}; the piecewise argument above is
given because the reconstruction of Sec.~\ref{sec:framework:kernel} is not
smooth, and because it is the argument the released code implements.

\subsection{The ladder is calibration free}
\label{sec:M:calibration}

The practical content of Eq.~\eqref{eq:rungs} is in the ratios. For any two
rungs,
\begin{equation}
  \frac{I_k}{I_j} = \frac{a_j}{a_k} ,
  \label{eq:ratio}
\end{equation}
in which $\Gps$, $\gamma$ and every response parameter cancel identically.

Reading the $a_k$ off a measurement needs one condition beyond (M), which
constrains the wavelength dependence and says nothing about the behavior of the
response near zero power:
\begin{quote}
\textbf{(L)} In the low-excitation limit the detected rate is linear in the
pumping rate with no offset, $R(\Gp) = R'(0)\,\Gp + o(\Gp)$ with $R'(0) > 0$,
after subtraction of background and dark counts.
\end{quote}
Hypothesis (L) is what ``unsaturated'' means operationally, and it is
independent of (M). Under (M) and (L) together, $R \propto A$ at low power, so
the critical values $a_k$ are read directly off a low-power photoluminescence
scan.

That scan must be taken \emph{at fixed incident optical power}, the convention
of Sec.~\ref{sec:framework:A}. A scan holding the incident photon flux
fixed instead returns $\sabs$ rather than $A = \lambda\sabs$; its critical
values are those of a different function, and the rung ratios of
Eq.~\eqref{eq:ratio} are then wrong by $\lambda_k/\lambda_j$, which reaches
$29\%$ across the analysis window---an order of magnitude beyond the
uncertainty the $a_k$ otherwise carry.

The positions of the
rungs are therefore predicted from a \emph{measured relative spectrum alone},
with no absolute power calibration and no knowledge of which mechanism
operates. Only the overall scale $I_c$ has to be supplied, and a
single-wavelength power sweep supplies it (test~T4 of Sec.~\ref{sec:tests}).

This is the sense in which Theorem~M functions as a prediction rather than a
fit. Nothing is adjusted between the low-power scan and the observed step
positions.

\subsection{Worked example at \SI{120}{\giga\pascal}}
\label{sec:M:example}

\begin{figure*}[t]
  \includegraphics[width=\textwidth]{figures/fig_m_level_set_ladder.pdf}
  \caption{The level set of Theorem~M and its multiplicity ladder, for the
  reconstructed kernel at \SI{120}{\giga\pascal}. (a) The normalized
  fixed-power figure of merit $a(\lambda) = A(\lambda)/A_{\max}$ with
  $A = \lambda\sabs$, dimensionless, against wavelength in
  \si{\nano\meter} over the analysis window
  \SIrange{402}{517}{\nano\meter}. Triangles mark the four interior maxima
  (Table~\ref{tab:maxima}); the dashed horizontal is one level $I_c/I$, and the
  circles where it cuts the curve are the members of the optimal set
  $\Lset(I)$ of Eq.~\eqref{eq:levelset}, which are exactly degenerate in
  $\eta$. (b) The resulting multiplicity $N = |\Lset|$, dimensionless, against
  $\log_{10}(I/I_c)$. Steps occur at the powers predicted by
  Eq.~\eqref{eq:rungs} from the critical values alone; the dashed verticals
  are those predicted positions, and the numeral above each is the rung index
  $k$ of Table~\ref{tab:ladder}, which gives the power ratio, the generating
  wavelength and the step for each; rungs 2 and 3 lie $0.3\%$ apart in power
  and share one numeral. Rung 4 is drawn in gray because it is generated by
  the edge of the analysis window rather than by a critical point of $a$. The $N=6$ plateau is only $\times1.003$ wide
  in power and is declared untestable. Both panels take the zero-phonon height
  at face value; correcting it against the published Debye-Waller factor
  (Appendix~\ref{sec:kernel:internal}, K3) moves that first step from
  $I/I_c = 1.44$ to $1.28$ and leaves every multiplicity unchanged.}
  \label{fig:levelset}
\end{figure*}

\begin{table}[tb]
  \caption{The multiplicity ladder of the reconstructed kernel at
  \SI{120}{\giga\pascal}, enumerated over $1 < I/I_c \le 6$. The index $k$
  numbers the rungs in order of power, and is the numeral each vertical line
  carries in Fig.~\ref{fig:levelset}(b); rungs 2 and 3 lie $0.3\%$ apart in
  power and share one numeral there. ``Observed''
  counts the members of $\Lset(I)$ found by direct enumeration; ``predicted''
  is Eq.~\eqref{eq:rungs} evaluated from the critical values of $a$ alone. The
  six agree to better than $0.01\%$, which checks the enumeration rather than
  the values; those carry about $\pm2\%$ from the $1\%$ height reconstruction of
  Appendix~\ref{sec:kernel:cross}, and their fourth and fifth digits are set by
  the \SI{0.005}{\nano\meter} critical-point grid. The listing is truncated at
  $I/I_c = 6$
  and further rungs exist above it. The first row inherits the zero-phonon
  uncertainty of Table~\ref{tab:maxima} and may lie anywhere in
  $1.28 \le I/I_c \le 1.44$; it remains the lowest rung across that range, so
  the $N$ and $\Delta N$ columns are unchanged by it. The values are conditional on the
  reconstruction; what Theorem~M freezes is the ladder \emph{structure}.}
  \label{tab:ladder}
  \begin{ruledtabular}
  \begin{tabular}{ccccccc}
    $k$ & $I_k/I_c$ obs. & $I_k/I_c$ pred. & from (nm) & kind & $N$ & $\Delta N$ \\
    \colrule
    1 & 1.44 & 1.44 & 514.5 & maximum   & $2\to4$ & $+2$ \\
    2 & 1.51 & 1.51 & 475.6 & maximum   & $4\to6$ & $+2$ \\
    3 & 1.52 & 1.52 & 474.5 & minimum   & $6\to4$ & $-2$ \\
    4 & 1.86 & 1.86 & 402.0 & edge\footnotemark[1] & $4\to3$ & $-1$ \\
    5 & 3.61 & 3.61 & 500.2 & maximum   & $3\to5$ & $+2$ \\
    6 & 4.30 & 4.30 & 497.9 & minimum   & $5\to3$ & $-2$ \\
  \end{tabular}
  \footnotetext[1]{Not a critical point of $A$. This rung is generated by the
  blue edge of the analysis window, which is where the source figure stops
  rather than where the absorption band ends
  (Appendix~\ref{sec:kernel:internal}, K1). Theorem~M's $\Delta N = -1$ edge
  case is general---any finite scan window produces one---but the position
  \SI{402}{\nano\meter} is a property of the reconstruction, not of the
  defect. Figure~\ref{fig:levelset}(b) draws this rung in gray.}
  \end{ruledtabular}
\end{table}

Table~\ref{tab:ladder} and Fig.~\ref{fig:levelset} give the ladder that the
reconstructed kernel produces at \SI{120}{\giga\pascal}, enumerated over
$1 < I/I_c \le 6$. Taking the zero-phonon height at face value, the
multiplicity runs $N = 2 \to 4 \to 6 \to 4 \to 3 \to 5 \to 3$, of which the
fourth step is the window-censoring event discussed below rather than a
critical point of $A$. All six
predicted step positions are reproduced to better than $0.01\%$, which checks
the level-set enumeration against Eq.~\eqref{eq:rungs} rather than the values
of $a$. The enumeration is truncated, not
complete: the count of Eq.~\eqref{eq:deltaN} requires a third interior
minimum at the foot of the
zero-phonon spike, and the red edge of the window supplies a further
$\Delta N = -1$, both of which fall above $I/I_c = 6$. The ladder continues.

\emph{The zero-phonon rung is the least certain in the table, and it is the one
that decides where the ladder starts.} The first step, at $I/I_c = 1.44$, is
generated
by the zero-phonon line (ZPL) maximum at \SI{514.5}{\nano\meter}. This rung, alone among the six,
depends on a zero-phonon \emph{height} rather than a position. Check~K3 of
Appendix~\ref{sec:kernel:internal} compares the kernel's own zero-phonon area
with the published Debye-Waller factor and finds it low by $11\%$, by the same
amount at every pressure. Correcting $a = 0.694$ by that constant gives
$a \simeq 0.780$ and moves the rung from $I/I_c = 1.44$ to $1.28$. The
defensible statement is therefore that the ZPL generates a rung somewhere in
$1.28 \le I/I_c \le 1.44$, and that its optical-limit sensitivity penalty lies
between $1.13$ and $1.20$. Neither end of that range is a bound: K3 measures
an offset in \emph{area}, while check~K2 finds the zero-phonon \emph{width} in
this reconstruction to be resolution limited, and an area ratio transfers to a
height only when the width is right. The range indicates a scale, not a limit.

That correction displaces one row of the ladder and does not reorder it. The
critical values of $a$ are $0.694$ at the zero-phonon maximum and then
$0.660$, $0.658$, $0.536$, $0.277$ and $0.233$; moving the first to $0.780$
carries it further from the rest rather than past any of them, so the
zero-phonon step remains the lowest rung and the sequence
$N = 2 \to 4 \to 6 \to 4 \to 3 \to 5 \to 3$ holds at both ends of the range.
What the uncertainty moves is where the ladder starts, not what it is. Test~T5
locates that first rung, and does so without absolute calibration.

What survives that uncertainty is the qualitative point, and it is the one that
matters for experimental design: a candidate-line list that omits the ZPL
cannot see this rung at any power. We are not aware of a candidate-line list
for high-pressure work that includes the zero-phonon line, whose
\emph{photoluminescence} penalty there is unremarkable. Its sensitivity penalty is
$1.13$ at the optimistic end of the range above and $1.20$ at the pessimistic
end, where it is indistinguishable from the $1.21$ of the commercial
\SI{473}{\nano\meter} line. On either reading it belongs on the list.

\emph{One plateau must not be claimed.} The plateau between the two rungs near
\SI{475}{\nano\meter} spans only $\times1.003$ in power. No realistic power
resolution reaches it, and we declare it untestable rather than report it. It
survives the zero-phonon uncertainty, since the pair of rungs that makes it
narrow does not involve the zero-phonon line and keeps its label $N = 6$ at
both ends of the range. The resolvable plateaux within the enumerated
range are, at face value, the initial $N = 2$ plateau of width
$\times1.44$, followed by $N = 4$,
$4$, $3$ and $5$ of width $\times1.05$, $\times1.23$, $\times1.93$ and
$\times1.19$. Correcting the zero-phonon height narrows the initial plateau to
$\times1.28$ and widens the first $N = 4$ plateau to $\times1.18$; the other
three are untouched, since no rung bounding them involves the zero-phonon
line. The widest plateau is the same on either reading---the $N = 3$ plateau
ending at the \SI{500.2}{\nano\meter} rung, at $\times1.93$---and every
resolvable plateau is wide enough to survive the correction.

\emph{The blue edge is a property of the figure, not of the defect.} The step
at $I/I_c = 1.86$ is generated by the edge of the analysis window at
\SI{402}{\nano\meter}, which is where the published spectrum stops rather than
where the absorption band ends (Appendix~\ref{sec:kernel}). It is counted as
$\Delta N = -1$ for that reason.

\subsection{The ladder is reachable}
\label{sec:M:reachability}

The ladder does not require exotic powers.
Reference~\cite{Dreau2011} reports closed forms for $R$, $C$ and
$\Delta\nu$ against optical pumping power, fitted to single-NV measurements;
Sec.~\ref{sec:G:measured} converts them into $E = 3$ with $n = 2$ and
$\rho^* = 1/5$, from which the sensitivity-optimal power is
$\varsigma^* = 0.072$ in units of the saturation power, for the representative
microwave drive $\Omega_R = \SI{3e6}{\radian\per\second}$. The lowest rung of
the \SI{120}{\giga\pascal} example then sits at
\begin{equation}
  \varsigma = 1.44 \times 0.072 = 0.10 ,
  \label{eq:reach}
\end{equation}
that is, at about a tenth of the saturation power---\emph{below} where many
experiments routinely operate rather than above it.

Two caveats attach to that figure, and neither removes the conclusion.
$\varsigma^*$ scales with the microwave drive (Sec.~\ref{sec:G:measured}), so
across the five settings of Ref.~\cite{Dreau2011} it runs from $0.025$ to
$0.140$. The zero-phonon uncertainty enters only weakly: the first rung moves
between $I/I_c = 1.28$ and $1.44$ and remains the lowest either way.
Combining both, the first rung falls between $3\%$ and $20\%$ of the saturation
power. Sweeping to ten
times saturation reaches $I/I_c \le 139$ at this drive, well beyond the highest
rung enumerated here. The ladder lies inside the accessible range by a
comfortable margin.

\subsection{The ladder is invariant across the degeneracy}
\label{sec:M:invariance}

The gauge degeneracy of Sec.~\ref{sec:G:identifiability}, which obstructs
mechanism attribution, protects the ladder.
Theorem~M's transition powers depend on the response only through $\Gps$, and
expressing them as $I_k/I_c$ removes $\Gps$ entirely. Hence the ladder is
invariant across the whole gauge plane. We check the collapse numerically
rather than assume it: giving the five models of Fig.~\ref{fig:gauge}(a) five
different half-scales $\Gamma$ gives them five different $\Gps$, and hence five
sets of \emph{absolute} probe intensities spanning a factor $16$. Probing
$N(I)$ at six points of each set---at $I/I_c = 1.05$, $1.45$, $1.55$, $2.0$,
$3.7$ and $4.5$ in each model's own units---returns the same multiplicities in
all five cases, namely $2$, $4$, $4$, $3$, $5$, $3$. It returns the same six at
both ends of the zero-phonon range of Sec.~\ref{sec:M:example}, because that
range moves the first rung between $1.28$ and $1.44$ and the first probe point
sits below both. The probe deliberately steps
over the unresolvable narrow window between the two rungs near
\SI{475}{\nano\meter}.

The ladder can therefore be predicted and tested \emph{without knowing which
mechanism operates}. Non-identifiability is an obstacle for mechanism
assignment and a robustness property for the central prediction of this work.

% =====================================================================
\section{Pressure-driven branch exchange}
\label{sec:X}

Theorem~M asks what power does to the optimum at one pressure. This section
asks the orthogonal question, and finds an effect the ladder cannot produce.

An absorption spectrum built on a Franck-Condon progression carries two
structurally distinct branches: the zero-phonon line and the phonon sideband.
Their weights scale differently with the Huang-Rhys factor $S$, the ZPL weight
collapsing as $e^{-S}$ while the sideband weight saturates as $1-e^{-S}$.

\begin{quote}
\textbf{Theorem X} (laser-dominated regime). Let $A_\ZPL(P)$ and $A_\SB(P)$ be
the fixed-power figures of merit of the two branches, and suppose the
excitation is laser dominated over the pressure range considered,
$W_{\mathrm{laser}} \ge \Gamma_\ZPL(P)$, so that
$\Weff = W_{\mathrm{laser}}$ is independent of $P$. Let
$D(P;\Weff) = \ln[A_\SB/A_\ZPL] = \ln[r(P)\,\Weff]$ with $r$ given by
Eq.~\eqref{eq:r}. If $r$ is strictly monotone in $P$, then $D$ has at most one
zero; and where $D$ changes sign the global optimum exchanges branch, the
optimal wavelength jumping from $\lambda_\ZPL(P^*)$ to $\lambda_\SB(P^*)$
\emph{taking no intermediate value}.
\end{quote}

\emph{Proof.} The two branch maxima are isolated, and between them $A$ is
strictly smaller than at either, so at every pressure the global maximizer of
$A$ lies in the two-point set $\{\lambda_\ZPL(P), \lambda_\SB(P)\}$. It is
$\lambda_\ZPL$ where $D<0$ and $\lambda_\SB$ where $D>0$, by the definition of
$D$. A map into a two-point set cannot take a value between the two points, so
at a sign change of $D$ the optimum moves by the full separation
$\lambda_\SB(P^*)-\lambda_\ZPL(P^*)$ rather than drifting through it. Monotone
$r$ gives $D$ at most one zero, hence at most one such jump. $\blacksquare$

The regime hypothesis is what carries the monotonicity, and it is worth being
explicit about why. Once $W_{\mathrm{laser}}$ falls below $\Gamma_\ZPL(P)$,
Eq.~\eqref{eq:weff} pins $\Weff$ to the line itself, $D$ becomes
$\ln[r(P)\,\Gamma_\ZPL(P)]$, and monotone $r$ no longer bounds the number of
zeros: that now requires the \emph{product} $r\,\Gamma_\ZPL$ to be monotone,
which the reconstruction cannot decide, because check~K2 of
Appendix~\ref{sec:kernel:internal} shows that $\Gamma_\ZPL$ is not readable
from it. Restricting the theorem is therefore not a caution but the condition
under which its proof runs.

\begin{quote}
\textbf{Corollary X$'$} (line-limited regime). In the complementary regime
$W_{\mathrm{laser}} < \Gamma_\ZPL(P)$, $D$ is independent of
$W_{\mathrm{laser}}$, and so is any crossing it has. Hence $P^*$ does not move
as the laser is narrowed further: the measured curve $P^*(W_{\mathrm{laser}})$
falls while the excitation is laser dominated and is flat below, and the knee
between the two regimes sits at $W_{\mathrm{laser}} = \Gamma_\ZPL(P^*)$.
\end{quote}

Corollary~X$'$ is an experimental extension rather than a second theorem: it
inherits Theorem~X's two-branch structure but not its monotonicity, and it
predicts the \emph{shape} of a curve rather than the existence of a crossing.
Its value is that the knee measures a quantity nothing else here can reach.

The theorem is stated in the bandwidth-resolved form because that is the form
this source can support. Writing the Franck-Condon pair as
$A_\ZPL \sim e^{-S}/\Gamma_\ZPL$ and $A_\SB \sim (1-e^{-S})/\Gamma_\SB$ gives
\begin{equation}
  \frac{\mathrm{d}D}{\mathrm{d}P}
  = \frac{1}{1-e^{-S}}\,\frac{\mathrm{d}S}{\mathrm{d}P}
  + \frac{\mathrm{d}}{\mathrm{d}P}\ln\frac{\Gamma_\ZPL}{\Gamma_\SB} ,
  \label{eq:dDdP}
\end{equation}
whose first term is positive whenever compression strengthens the
electron-phonon coupling, since $1/(1-e^{-S}) \ge 1$. That would be sufficient
for monotonicity if the ZPL also broadened at least as fast as the sideband,
but Appendix~\ref{sec:kernel:internal} shows the zero-phonon width is not readable
from this source at all, so we do not use it. That is precisely why Theorem~X
is stated in the laser-dominated regime: there $\Weff = W_{\mathrm{laser}}$ is
independent of pressure, $\Gamma_\ZPL$ leaves the problem, and monotonicity
rests on $r(P)$ alone. Section~\ref{sec:X:bandwidth} establishes it there, from
published data. In the line-limited regime of Corollary~X$'$ the width does not
leave the problem, and the missing $\Gamma_\ZPL$ is what stops us asserting
monotonicity there.

Both antecedents are generic for a color center under compression, so what
varies between materials is \emph{where} $P^*$ falls inside the accessible
range rather than \emph{whether} an exchange occurs. That is the general
condition the framework offers to high-pressure optical sensing: a once-only
exchange, generic rather than an accident of one host.

The jump is large enough to matter to an experiment. In the reconstructed
kernel the two branch maxima are separated by
\SIrange{68}{78}{\nano\meter} at every published pressure---at
\SI{120}{\giga\pascal} they sit at \SI{514.5}{\nano\meter} and
\SI{440.6}{\nano\meter}, \SI{73.9}{\nano\meter} apart---so wherever $P^*$ falls
in the accessible range, the recommended excitation line moves by roughly
\SI{70}{\nano\meter} at that pressure instead of drifting toward it. No single
line is correct on both sides.

\subsection{The drivers, from published data}
\label{sec:X:antecedents}

The electron-phonon driver is read from the tabulated Huang-Rhys and
Debye-Waller factors of Ref.~\cite{Ho2026} itself, independently of our
reconstruction of its spectra. The absorption Huang-Rhys factor rises from
$3.081$ to $4.613$ over
\SIrange{0}{120}{\giga\pascal}, an increase of $50\%$, monotone, at
$\mathrm{d}\Sabs/\mathrm{d}P = \SI{1.27e-2}{\per\giga\pascal}$. The absorption
Debye-Waller factor falls from $2.180\times10^{-2}$ to $3.63\times10^{-3}$, a
factor $6.01$, monotone; both endpoints are the values Ref.~\cite{Ho2026}
states in its own text, and it describes the fall as more than fivefold. The
zero-phonon weight therefore collapses in the published data more steeply than
a single effective mode gives: $e^{-S}$ with the tabulated $S$ would produce
only a factor $4.63$, so the published Debye-Waller factor is multi-mode,
though by a margin of $1.30$ rather than a wide one.

The second term of Eq.~\eqref{eq:dDdP}, the relative broadening of the two
branches, is \emph{not} confirmed and cannot be from this source: the apparent
zero-phonon width in the published spectra is resolution limited
(Appendix~\ref{sec:kernel:internal}, K2). Section~\ref{sec:X:bandwidth}
therefore works in the laser-dominated regime, where the effective bandwidth
replaces that width, and establishes the monotonicity of the ratio empirically
instead.

Three checks support the two-branch identification. First, the two branches
appear as independent local maxima of the \emph{raw extracted samples}---not of
any interpolation---at all seven published pressures. The zero-phonon branch at
ambient pressure lands at \SI{637.1}{\nano\meter}, i.e.
\SI{1.946}{\electronvolt}, which is the \NVm{} zero-phonon line and is a value
fitted nowhere in the pipeline. And the two branches carry different pressure
coefficients, $3.8 \pm 0.2$ against
$5.1 \pm 0.2$~\si{\milli\electronvolt\per\giga\pascal} from the seven
published positions, so they are not one feature counted twice;
correspondingly the Franck-Condon displacement $S\hbar\omega$ grows from
\SIrange{0.232}{0.404}{\electronvolt}, an increase of $74\%$. Both
coefficients come from one reconstruction of one source; the independent
high-pressure optical measurements of Ref.~\cite{Lyapin2018} offer the
external comparison that would make them harder to attribute to the
reconstruction, and we have not made it. The three checks together identify the
two branches as distinct features of the source rather than one feature
resolved twice.

\subsection{The exchange pressure depends on the excitation bandwidth}
\label{sec:X:bandwidth}

\begin{figure*}[t]
  \includegraphics[width=\textwidth]{figures/fig_x_branch_exchange.pdf}
  \caption{The branch exchange of Theorem~X. SB, phonon sideband; ZPL,
  zero-phonon line. (a) The branch ratio per unit excitation bandwidth,
  $r(P) = \lambda_\SB\sabs^\SB/(\lambda_\ZPL\DWF_{\mathrm{abs}})$ of
  Eq.~\eqref{eq:r}, in \si{\per\electronvolt}, against pressure in
  \si{\giga\pascal} on a log ordinate. It rises monotonically by a factor
  $4.3$ over \SIrange{0}{120}{\giga\pascal}, which is what makes the crossing
  of Theorem~X unique where it occurs. The right axis is the critical
  bandwidth $1/r$ in \si{\milli\electronvolt}, below which the zero-phonon
  branch is competitive; the annotated endpoints are
  \SIlist{10.30;2.39}{\milli\electronvolt}. (b) The exchange pressure $P^*$,
  in \si{\giga\pascal}, against the effective excitation bandwidth $\Weff$
  of Eq.~\eqref{eq:weff}, in
  \si{\milli\electronvolt}, at fixed power. Square markers are six
  bandwidths spanning the accessible window,
  \SIlist{2.5;3.0;4.0;5.0;7.0;9.0}{\milli\electronvolt}; the numeral beside
  each is the exchange pressure it produces, and $P^*$ carries
  $\pm\SI{2}{\giga\pascal}$. The two dashed verticals bound that window and
  the shaded bands are outside the published pressure range,
  which no crossing enters for $\Weff$ outside roughly
  \SIrange{2.4}{10.3}{\milli\electronvolt}. Sweeping the laser linewidth moves
  $P^*$ monotonically, and no other mechanism in this framework does so.}
  \label{fig:branch}
\end{figure*}

The two branches cannot simply have their peak heights compared, because they
are different kinds of object. The sideband enters through a lineshape
\emph{density} (per unit energy), while the ZPL enters through a
\emph{dimensionless weight} that must be divided by whatever bandwidth samples
it. That sampling width is not the laser linewidth alone. The absorbed
zero-phonon rate is set by the overlap of the laser spectral profile with the
intrinsic zero-phonon lineshape, so it grows as the laser is narrowed only
until the laser is narrower than the line itself, after which it saturates. At
the level of widths,
\begin{equation}
  \Weff \simeq \max\!\left(W_{\mathrm{laser}},\,\Gamma_\ZPL(P)\right) ,
  \label{eq:weff}
\end{equation}
and it is $\Weff$, not $W_{\mathrm{laser}}$, that enters. Writing
\begin{equation}
  r(P) = \frac{\lambda_\SB\,\sigma_\SB}
              {\lambda_\ZPL\,\DWF_{\mathrm{abs}}}
  \quad [\mathrm{eV}^{-1}] ,
  \qquad
  \frac{A_\SB}{A_\ZPL} = r(P)\,\Weff ,
  \label{eq:r}
\end{equation}
the exchange occurs where $r(P) = 1/\Weff$. Table~\ref{tab:drivers} shows that $r$
is monotone increasing, by a factor $4.31$ over the published range, so at most
one crossing survives and Theorem~X's monotonicity antecedent holds across that
range. What changes is that $P^*$ is no
longer a single number, and that for bandwidths outside
\SIrange{2.4}{10.3}{\milli\electronvolt} no crossing falls inside the accessible
range at all.

\begin{table}[tb]
  \caption{The drivers of Theorem~X. DWF, Debye-Waller factor; SB, phonon
  sideband. Taken from the published theory of
  Ref.~\cite{Ho2026} [its panels (b) and (c)] rather than from the
  reconstructed spectra. $r(P)$ is the branch ratio per unit excitation
  bandwidth, Eq.~\eqref{eq:r}; $1/r$ is the bandwidth below which the ZPL
  branch is competitive.}
  \label{tab:drivers}
  \begin{ruledtabular}
  \begin{tabular}{lccccccc}
    $P$ (GPa) & 0 & 20 & 40 & 60 & 80 & 100 & 120 \\
    \colrule
    $\Sabs$ & 3.081 & 3.368 & 3.648 & 3.908 & 4.152 & 4.375 & 4.613 \\
    $\DWF_{\mathrm{abs}}$ ($10^{-2}$)
      & 2.180 & 1.565 & 1.127 & 0.835 & 0.620 & 0.479 & 0.363 \\
    $r$ (eV$^{-1}$)
      & 97.1 & 121.9 & 162.1 & 208.0 & 267.1 & 329.1 & 418.4 \\
    $1/r$ (meV)
      & 10.30 & 8.20 & 6.17 & 4.81 & 3.74 & 3.04 & 2.39 \\
  \end{tabular}
  \end{ruledtabular}
\end{table}

Figure~\ref{fig:branch}(b) gives the dependence, at fixed incident optical
power: $P^*$ falls from \SI{117}{\giga\pascal} at
$\Weff = \SI{2.5}{\milli\electronvolt}$ through \SI{101}{\giga\pascal} at
\SI{3}{\milli\electronvolt} to \SI{13}{\giga\pascal} at
\SI{9}{\milli\electronvolt}, each value carrying about
\SI{2}{\giga\pascal} from the $3\%$ uncertainty on $r$; at fixed incident
photon flux each moves by about \SI{10}{\giga\pascal}. A laser linewidth
chosen anywhere in the commercial range therefore moves the exchange across
the whole accessible pressure axis. Equivalently, the ZPL competes at all
only for an excitation bandwidth below $1/r(P)$, which narrows from
\SI{10.3}{\milli\electronvolt} at ambient pressure to
\SI{2.4}{\milli\electronvolt} at \SI{120}{\giga\pascal}.

The bandwidth dependence is a sharply testable prediction, and Eq.~\eqref{eq:weff}
sharpens it further. By Theorem~X the exchange pressure moves monotonically
with the laser linewidth while the excitation is laser dominated, and by
Corollary~X$'$ it stops moving below that. A measured
$P^*(W_{\mathrm{laser}})$ curve should therefore fall and then flatten, and the
\emph{knee} between the two regimes locates $\Gamma_\ZPL$ at that
pressure---a quantity check~K2 of Appendix~\ref{sec:kernel:internal} shows
cannot be read from the reconstruction at all. No other mechanism in this
framework produces either the falling segment or the knee; a fixed-wavelength
ranking crossover has neither. It also raises the value
of test~T11 of Sec.~\ref{sec:tests}, which probes the electron-phonon driver
without touching the bandwidth question at all.

\emph{Reporting rule.} A value of $P^*$ must never be quoted without stating
both the excitation bandwidth and the fixed-power versus fixed-flux convention.

\subsection{What this is not}
\label{sec:X:not}

Two nearby phenomena would, if conflated with branch exchange, make a
generic effect look like an accident.

\emph{It is not a fixed-wavelength ranking crossover.} In the same kernel the
sensitivity ranking of \SIlist{457;532}{\nano\meter} reverses near
\SI{51}{\giga\pascal}, which is the geometry of a single band sweeping past the
midpoint of two fixed probes and requires no branch structure at all. That
pressure corresponds to a bandwidth of about \SI{5}{\milli\electronvolt} and
lies inside the branch-exchange band of Fig.~\ref{fig:branch}(b), so the two
cannot be told apart by where they fall. The discriminant is mechanism, and it
is test~T9: a ranking crossover involves one peak, an exchange requires
\emph{both} branches on either side of $P^*$.

\emph{It is not visible to a single-mode envelope.} A single-effective-mode
Franck-Condon envelope has exactly one interior local maximum at every pressure
between $0$ and \SI{120}{\giga\pascal}, so it has no branches to exchange and
any crossover it shows is the fixed-wavelength kind, at a pressure of its own.
The reliable discriminant is to count the interior local maxima of $A(\lambda)$
in each model---the same count that decides whether Theorem~M has a ladder at
all---rather than to compare two pressures and read a coincidence as agreement.

\subsection{The \texorpdfstring{$(P,I)$}{(P,I)} plane}
\label{sec:phase}

Theorems~M and~X produce degeneracies of different origin---one along the power
axis, one along the pressure axis---and they stack.

\begin{itemize}
  \item \emph{Below $P^*$}, the zero-phonon branch carries the global optimum,
  so the first members of $\Lset(I)$ lie on it. The level set is defined over
  the whole window, so as the level falls it crosses critical values on the
  sideband as well, and those wavelengths join the optimal set.
  \item \emph{At $P^*$}, the low-power optimum is already twofold, one member
  on each branch, before any power-induced structure appears. Raising the power
  then adds rungs drawn from both. This is the point of richest structure.
  \item \emph{Above $P^*$}, the sideband branch carries the global optimum and
  supplies the first members; the zero-phonon branch enters the optimal set
  once the level falls to its peak.
\end{itemize}

The ladder therefore mixes branches rather than being confined to one, and the
worked example shows it doing so. At \SI{120}{\giga\pascal}, which lies at or
above $P^*$ for every bandwidth in Fig.~\ref{fig:branch}(b), one rung of
Table~\ref{tab:ladder} is generated by the zero-phonon maximum, four are
sideband features and one is the window edge---and on the face-value branch the
zero-phonon rung is the lowest of them. What the pressure fixes is which branch
carries
the \emph{global} optimum, and hence the order in which the two branches enter
the ladder; it does not assign the ladder to a branch.

Theorem~M is thus the structure along the $I$ axis and Theorem~X the structure
along the $P$ axis; together they give the phase diagram of the optimal
operating point.

The practical discriminant between them is that \emph{Theorem~X's degeneracy
exists at zero power while Theorem~M's requires $I > I_c$}. At $P^*$ the
optimum is already
twofold in the low-excitation limit, so the two can be separated by measuring
at low power and varying pressure, before any power-induced multiplicity
appears. Test~T10 is that separation.

% =====================================================================
\section{Pre-specified tests}
\label{sec:tests}

Each of the three theorems makes predictions that can fail, and the value of
stating them in advance is that none can be adjusted after the data are in.
Table~\ref{tab:tests} collects the eleven tests they license, grouped by the
claim each one probes.

\begin{table*}[t]
  \caption{Pre-specified tests. Each entry states what an outcome is
  consistent with and what it falsifies. Agreement of finitely many observables
  at finite precision constrains (M) and does not establish it: a
  photon-energy-specific channel that happened to track $A(\lambda)$ over the
  measured range would pass every forward test. (M) is therefore strongly
  falsifiable and not uniquely provable from this set.
  T1--T4 probe the mediation hypothesis (M) and
  the coincidence condition; T5--T7 test Theorem~M and Theorem~G; T8--T11 test
  Theorem~X. None requires absolute calibration: every entry is a ratio, a step
  position, or an equality.}
  \label{tab:tests}
  \newcommand{\Pc}[1]{\parbox[t]{0.325\textwidth}{\raggedright #1\strut}}
  \newcommand{\Dc}[1]{\parbox[t]{0.395\textwidth}{\raggedright #1\strut}}
  \begin{ruledtabular}
  \begin{tabular}{lll}
    test & procedure & discriminates \\
    \colrule
    T1 collapse
      & \Pc{plot $\eta$ against the measured $R$ for all $(\lambda,I)$,
           discarding $\lambda$}
      & \Dc{collapse is consistent with (M); the scatter
           measures $\ell_G$} \\
    T2 power closure
      & \Pc{measure $\leta(I) - \lPL(I)$ over a power sweep}
      & \Dc{closes as $I\to0$: consistent with intensity mediation; a
           persistent offset falsifies (M)} \\
    T3 degeneracy
      & \Pc{compare $\eta$ at the members of a doublet, searching for equal
           $\eta$ rather than symmetric placement}
      & \Dc{equality within the measurement error supports mediation by the
           pump rate; an imbalance measures $\ell_G$ and falsifies (M)} \\
    T4 single-line prediction
      & \Pc{from one power sweep at a single wavelength, extract the response
           scales and hence $\Gps$}
      & \Dc{predicts $I_c$ and the ladder \emph{without any wavelength scan}} \\
    \colrule
    T5 ladder
      & \Pc{read the critical values $a_k$ off a low-power $R(\lambda)$ scan
           at fixed incident optical power;
           predict $I_k/I_c = 1/a_k$; verify the steps}
      & \Dc{steps at the predicted powers support Theorem~M; steps
           elsewhere, or none, falsify it} \\
    T6 universal degeneracy
      & \Pc{compare $\eta$ across \emph{all} $N$ members of a multiplet}
      & \Dc{equality at the noise floor bounds any violation; the
           deviation measures it} \\
    T7 exponent closure
      & \Pc{measure $C(I)$, $R(I)$, $\Delta\nu(I)$ separately; form
           $E = 2c+s+2w$ and the pump nonlinearity $n$; predict
           $\rho^* = 1/(En-1)$}
      & \Dc{agreement is consistent with a power-law response; disagreement
           falsifies the assumed functional form} \\
    \colrule
    T8 exchange
      & \Pc{track $\lPL(P)$ at low power through $P^*$}
      & \Dc{a discontinuous jump supports Theorem~X; continuous motion
           falsifies it} \\
    T9 branch coexistence
      & \Pc{check that both peaks are present on either side of $P^*$}
      & \Dc{both present is consistent with exchange; only one indicates
           branch disappearance, a different mechanism} \\
    T10 zero-power degeneracy
      & \Pc{at $P \approx P^*$ compare $\eta$ at the two wavelengths in the
           low-excitation limit}
      & \Dc{equal at zero power indicates Theorem~X; a power threshold
           indicates Theorem~M} \\
    T11 driver
      & \Pc{measure the branch separation $S\hbar\omega$ against pressure, and
           invert $\psi(p^*+1) = \ln S$ to isolate $S$}
      & \Dc{monotone increase of $S$ establishes the antecedent of
           Theorem~X; a constant $S$ predicts no exchange}
  \end{tabular}
  \end{ruledtabular}
\end{table*}

\emph{T4 converts a single-wavelength measurement into a prediction about a
multi-wavelength scan.} On synthetic data---fourteen powers at one wavelength,
$2\%$ multiplicative noise---$\Gps$ is recovered to $2.7\%$ in a single trial,
and across $200$ trials the $16$--$84\%$ interval is about $\pm2\%$. Combined
with T5, which supplies the $a_k$ from a low-power scan, the whole ladder is
predicted before any wavelength-resolved sensitivity measurement is made.

\emph{T6 is a null test of arbitrary precision.} Under (M) the degeneracy among
the members of $\Lset$ is exact, so any measured deviation is itself a
measurement of how far (M) fails.

\emph{T11 is the cheapest.} The branch separation is a difference of two peak
positions, so it needs neither absolute nor intensity calibration. One caution:
the separation is $S\hbar\omega$, and compression hardens $\hbar\omega$ as well
as raising $S$. Isolating $S$ requires the stationarity condition of the
Franck-Condon envelope, $\psi(p^*+1) = \ln \Sabs$ with $\psi$ the digamma
function and $p^*$ the observed sideband maximum in units of $\hbar\omega$; the
continuum shortcut $p^* = \Sabs$ is wrong by about half a phonon over the
relevant range. If the separation does not grow with pressure, the
electron-phonon driver of Theorem~X fails and no exchange is predicted.

We also record the falsification conditions explicitly. If T1 collapses onto a
single curve \emph{and} $\leta = \lPL$ within \SI{3}{\nano\meter} at every
power, then no divergence exists, the conventional practice of maximizing
photoluminescence is correct, and the central proposition of this work fails.
If the predicted rungs $I_k/I_c$ carry no steps within the power resolution,
Theorem~M fails---either (M) does not hold or $\Phi$ has more than one interior
maximum. If $N(I)$ moves continuously rather than in steps, $\eta$ is not a
function of $\Gp$ alone and (M) is violated. If $\lPL(P)$ moves continuously
through $P^*$, or if the laser linewidth can be varied without moving $P^*$,
Theorem~X and its bandwidth corollary fail respectively.

% =====================================================================
\section{Discussion}
\label{sec:discussion}

Two of the framework's three antecedents have already been measured by
others, one has not, and the structures it reports have partial
precedents in sensing that are worth separating from it precisely. What
follows takes those in turn, then asks whether charge conversion alone could
produce the same effect, what the framework settles that was previously open,
why the effect belongs to high pressure, and where its limits lie.

\subsection{Relation to what has already been measured}
\label{sec:discussion:prior}

Two of the three antecedents this framework needs have already been measured.
The interior maximum of $\Phi$---the condition $En > 1$ without which there is
no level set, no degeneracy and no ladder---is the content of the closed forms
Ref.~\cite{Dreau2011} fitted to measurement in 2011, robust across all the
microwave settings that work reports (Sec.~\ref{sec:G:measured}). The
wavelength dependence of the ODMR contrast, $\ell_G \neq 0$, has been measured
at ambient pressure by Ref.~\cite{Todenhagen2023}. Neither is assumed here.

What that work leaves open is what remains. It optimizes the contrast rather
than the sensitivity, so the linewidth---which enters $\eta$ with the same
power as the contrast---is not recorded; and it re-optimizes the excitation
power at each wavelength, which is exactly the operation that destroys the
fixed-power comparison Theorem~M is built on, collapsing a two-dimensional
optimization over $(\lambda, I)$ onto a section through it. The multiplicity of
the optimal set, its degeneracy, the ladder and the gauge structure are
therefore untouched. One measurement would connect that work to this one:
contrast and linewidth recorded separately at a common power.

\subsection{Related notions of multiplicity and of exchange}
\label{sec:discussion:related}

Neither of the two structures this paper reports is without precedent in
sensing, and it is worth saying precisely where each departs from what exists.

Multiplicity first. Ma \emph{et al.} established that an absorption-based
sensor is designed around a \emph{set} of transitions rather than one, the set
being chosen to condition the inverse problem that maps measured absorbances
onto a nonuniform temperature field~\cite{Ma2010}. Sun \emph{et al.} treated
the several resonances of a nonlinear photonic gyroscope as a single joint
resource and maximized the Fisher information of the whole
cavity~\cite{Sun2024}. In the first the multiplicity is selected by the
experimenter and its members carry different information; in the second the
resonances cooperate, and the figure of merit is a functional of all of them
together. Here the members of $\Lset(I)$ are alternatives, each individually
optimal and mutually degenerate in $\eta$, nothing selects them, and it is
their \emph{number} rather than their combination that is the observable.

Exchange second. Shaji Rebeirro \emph{et al.} demonstrated microwave-free
wide-field NV magnetometry at the ground-state level anticrossing, one member
of a family of operating features that also contains the excited-state
anticrossing and the zero-field crossings~\cite{ShajiRebeirro2024}. Ronchi
\emph{et al.} showed that the quantum Fisher information of a probe at an
avoided crossing is maximized by coherent evolution in one region of the
offset-against-available-time plane and by repeated projective measurement in
the other, so that the optimal strategy exchanges across a boundary in that
plane~\cite{Ronchi2024}. The branches of the first are branches of the static
bias field and the operator selects among them by setting that field; what
exchanges in the second is the measurement protocol applied to one transition.
Here neither the branch nor the protocol is chosen: which physical transition
carries the optimum is fixed by the pressure, under a protocol held constant,
and the exchange is a discontinuity in an external control variable rather than
a decision.

\subsection{Charge conversion alone cannot split the optima}
\label{sec:coincidence:P7}

It is natural to expect the NV charge state to be what separates the two
optima, since it is the most wavelength-sensitive quantity in the problem. It
cannot be, and the reason is structural. In the frozen charge model both
ionization from the excited state and recombination through excitation of
\NVz{} scale with the same absorption cross section, so the steady-state \NVm{}
fraction
\begin{equation}
  f_- = \frac{r_0\sigma + r_{\mathrm{bg}}}
             {(r_0 + a_{\mathrm{es}})\sigma + r_{\mathrm{bg}}}
  \;\xrightarrow[\;\sigma \gg r_{\mathrm{bg}}/r_0\;]{}\;
  \frac{r_0}{r_0 + a_{\mathrm{es}}}
  \label{eq:fminus}
\end{equation}
is a constant wherever the cross section is appreciable. Here $\sigma \equiv
\sabs$, $r_0$ is the recombination-rate coefficient for excitation of \NVz{},
$a_{\mathrm{es}}$ the excited-state ionization coefficient, and
$r_{\mathrm{bg}}$ a cross-section-independent background recombination rate.
More generally, $f_-$ is mediated by $\Gp$ in the sense of (M), so it enters
the three cases of Sec.~\ref{sec:coincidence:P3} through $A$ alone and cannot
itself be the source of a separation between $\leta$ and $\lPL$.

The corollary is sharp and identifies the only charge-related channel that
survives: charge conversion can move the sensitivity optimum \emph{only} if
$\sabs^{\NVz}$ has a spectral shape different from that of $\sabs^{\NVm}$, so
that the two rates cease to be proportional. That is a violation of (M), not an
instance of it. It is also the channel that the \SI{575}{\nano\meter} feature
of Ref.~\cite{Todenhagen2023}---the \NVz{} zero-phonon line---is evidence for,
and the channel that Ref.~\cite{Beha2012} probes. The optimum excitation
wavelength reported there for single centers at ambient pressure is an optimum
of the \emph{rate}, arising from the recharging balance; if that balance
carries wavelength structure of its own, it is structure of exactly this kind.
What Eq.~\eqref{eq:fminus} shows is that the frozen-charge part of the balance
cannot supply it.


\subsection{An open question this framework answers}
\label{sec:discussion:open}

Reference~\cite{Barry2020}, reviewing absorption-based readout, records that
the anomalous contrast reveals a strong wavelength \emph{and} power dependence
which suggests that green absorption readout is hindered by an unknown effect,
that the wavelength and power dependence of that effect suggests \NVz{}/\NVm{}
charge dynamics could play a role, and that further investigation might reveal
presently unknown NV dynamics~\cite{Barry2020,Bauch2010}.

That is a description of the $(\lambda, I)$ plane this paper organizes, and
Sec.~\ref{sec:coincidence:P7} answers it: charge conversion alone cannot
produce the effect, and the only channel that can is a difference in spectral
\emph{shape} between the \NVz{} and \NVm{} bands.

\subsection{Why the effect belongs to high pressure}
\label{sec:discussion:pressure}

Theorem~M is a statement about any absorbing sensor and could be tested at
ambient pressure; only the values of its rungs are pressure specific. Theorem~X
could not: it needs $\mathrm{d}S/\mathrm{d}P > 0$ (Sec.~\ref{sec:intro}), so
whether $P^*$ falls inside the accessible range is what decides observability
in a given host.

\subsection{Limitations}
\label{sec:discussion:limits}

Five limitations remain.

\emph{(M) is unverified, and is known to fail at discrete resonances.} The
resonances identified so far~\cite{Todenhagen2023} lie redward of
\SI{500}{\nano\meter}, outside the blue window that carries the high-pressure
optimum, but they are ambient-pressure measurements and the levels shift under
compression. We therefore present (M) as the object of test~T6, and claim only
that it can be tested in the blue window to arbitrary precision.

\emph{Every number is conditional on one reconstruction}, with node spacing of
\SIrange{3}{10}{\nano\meter} in the blue window and piecewise-linear
interpolation between nodes; the zero-phonon critical \emph{value} carries in
addition the constant $11\%$ offset against the published Debye-Waller factor
measured by check~K3 of Appendix~\ref{sec:kernel:internal}, which places the
first rung to about $\pm6\%$ in power. The
theorems are what we freeze. In an experiment the critical values $a_k$ are
read from the measured spectrum, not taken from Table~\ref{tab:ladder}.

\emph{Theorem~M assumes a unique interior maximum of $\Phi$.} If there are
several, the optimal set is a union of level sets and the ladder is
correspondingly more complex. For the power-law responses of
Eq.~\eqref{eq:powerlaw} with $En>1$ the maximum is unique, so test~T7 also
constrains this premise.

\emph{Theorem~X assumes two branches.} A spectrum with three or more
structurally distinct branches would exchange more than once, giving a sequence
of exchanges rather than a single one---the pressure-axis analogue of the
multiplicity ladder.

\emph{ZPL excitation is bandwidth sensitive in a way the other lines are not,
and the sensitivity is shared by Theorems~M and~X.}
If the laser linewidth exceeds the ZPL width, the effective $A(\lambda_\ZPL)$
falls and the lowest rung of the ladder moves to higher power. Any experiment
including the ZPL among its candidate lines must record the laser linewidth,
for the same reason that $P^*$ must be quoted with a bandwidth.

% =====================================================================
\section{Conclusions}
\label{sec:summary}

Selecting the excitation wavelength of an NV magnetometer by maximizing its
photoluminescence requires $C/\Delta\nu$ to be wavelength independent. That
condition is testable, and it is already known to fail. Relaxing it does not
displace the optimum; it changes what kind of object the optimum is.

Under the single hypothesis (M), and with no assumption about the absorption
spectrum, the sensitivity optimum is a level set whose members are exactly
degenerate and whose cardinality is a staircase in $\ln I$. The response enters
that staircase only through $E = 2c+s+2w$ and the pump nonlinearity $n$, so the
mechanism cannot be recovered, within that response class, from sensitivity
data at any number of powers,
and the ladder is correspondingly robust to not knowing it. And because
compression raises the Huang-Rhys factor, the branch carrying the global
optimum is exchanged discontinuously, at a pressure the experimenter selects
through the excitation bandwidth.

Two of the three antecedents have already been measured by others, at ambient
pressure, and the lowest rung of the
ladder falls between $4\%$ and $21\%$ of saturation---inside the range at which
experiments already operate. Eleven pre-specified tests follow, none requiring
absolute calibration; T11 is the cheapest, needing only two peak positions per
pressure, and T6 is a null test of arbitrary precision for (M) itself. In that
sense the framework supplies a way to measure its central hypothesis instead of
assuming it.

The stakes are set by Eq.~\eqref{eq:etaB}. Map acquisition time scales as
$\eta^2$, so the commercial \SI{473}{\nano\meter} line costs a factor $1.45$ in
time against the optimum at \SI{120}{\giga\pascal}, conditional on the
reconstruction (Appendix~\ref{app:optical})---and on the far side of $P^*$ the
correct line lies on the other branch entirely.

% ACKNOWLEDGMENTS.  Deliberately absent from the typeset output rather than
% carrying a placeholder: an empty section is worse than no section, and the
% funding line is a disclosure that has to be accurate.  Restore before the
% final version with (i) the grant agency and award numbers, (ii) the facility
% or beamline if applicable, (iii) individual thanks.  APS accepts an
% acknowledgment added at the proof stage.
%
% \begin{acknowledgments}
%   We thank ... .  This work was supported by ... under Grant No. ... .
% \end{acknowledgments}

% =====================================================================
\appendix

\section{Fixed-power optical limit at \SI{120}{\giga\pascal}}
\label{app:optical}

The results of the main text are structural. For reference we record the
fixed-power optical limit of the reconstructed kernel, which is the input the
theorems act on and the quantity a laser specification would be drawn from.

Assumptions: fixed incident optical power, low saturation, and
wavelength-independent charge yield, ODMR contrast, linewidth and collection
within the scan---that is, $\ell_G = 0$, so that by Eq.~\eqref{eq:P1b} the two
optima coincide and $\eta \propto a^{-1/2}$. Under exactly these assumptions
the interpolant maximum is \SI{440.6}{\nano\meter}, reported as approximately
\SI{441}{\nano\meter} because the extracted blue-side nodes are
\SIrange{3}{10}{\nano\meter} apart. The $5\%$ penalty band is
\SIrange{426.4}{457.9}{\nano\meter} and the $10\%$ band
\SIrange{420}{465}{\nano\meter}. Table~\ref{tab:maxima} gives the
penalties.

Of the commercial lines conventionally proposed for high-pressure NV work,
\SI{445}{\nano\meter} and \SI{457}{\nano\meter} fall inside the $5\%$ band, at
penalties of $1.004$ and $1.045$; the \SI{473}{\nano\meter} line costs $21\%$. The ground-state photoionization
threshold at \SI{120}{\giga\pascal}, $\mathrm{IP}(\gsA) =
\SI{3.06}{\electronvolt}$ or \SI{405.2}{\nano\meter}, imposes a hard blue
boundary: it is active only below \SI{405}{\nano\meter} and therefore sets a
wall without displacing a \SI{440}{\nano\meter} optimum. It is one of the
identified channels through which (M) can fail, but it fails outside the window
of interest.

\subsection{Extrapolation to 300--400 GPa: a blue wall, not a unique optimum}
\label{app:extrapolation}

The preceding result is anchored at \SI{120}{\giga\pascal}. It is nevertheless
useful to ask what a controlled continuation would imply for the
\SIrange{300}{400}{\giga\pascal} regime. We use the pressure-dependent ZPL and
ground-state ionization energy in Ref.~\cite{Ho2026} as the anchor, and impose
the physically motivated constraint that they are continued by the same law:
the authors report that $\mathrm{IP}(\gsA)$ closely follows the ZPL shift.
For any such continuation, the one-photon ionization boundary is
\begin{equation}
  \lambda_{\mathrm{ion}}(P)
  = \frac{hc}{\mathrm{IP}(\gsA,P)},
  \qquad hc=1239.84\ {\rm eV\,nm} .
  \label{eq:ionization_wall}
\end{equation}
Excitation at $\lambda<\lambda_{\mathrm{ion}}$ supplies a photon above the
ground-state threshold and can deplete NV$^-$ into NV$^0$; the corresponding
photoionization channels and their state dependence are established by the
ab-initio analysis of Ref.~\cite{Razinkovas2021}.

Table~\ref{tab:extrapolation} gives the resulting estimates at the target
pressures. The optimum range is the envelope of the internally consistent
variants: saturating versus linear ZPL continuation and the two illustrative
phonon-hardening factors. The \SI{120}{\giga\pascal} row is deliberately not
replaced: its independently reconstructed kernel optimum is the
\SI{440.6}{\nano\meter} value reported above.
The pressure-dependent optical kernel itself is measured only to
\SI{120}{\giga\pascal}, so this envelope is an uncertainty statement, not a
prediction of a unique laser line.

\begin{table}[t]
\caption{Illustrative continuation above the \SI{120}{\giga\pascal} anchor.
The ionization wall uses the same continuation for the ZPL and
$\mathrm{IP}(\gsA)$. The optimum interval spans the consistent model variants;
it should not be read as a confidence interval.}
\label{tab:extrapolation}
\begin{ruledtabular}
\begin{tabular}{c c c c}
 $P$ (GPa) & $\mathrm{IP}(\gsA)$ (eV) & $\lambda_{\mathrm{ion}}$ (nm) &
 $\lambda_{\mathrm{opt}}$ (nm) \\
 \hline
 300 & 3.63 & 341.6 & 349--451 \\
 400 & 3.95 & 314.1 & 314--448 \\
\end{tabular}
\end{ruledtabular}
\end{table}

The practical conclusion is one-sided. Once the optimum reaches the blue wall,
making the excitation wavelength still shorter cannot improve the sensing
rate: it increases the photon energy beyond the ionization threshold and
reduces the NV$^-$ population that produces ODMR. In the linear, coupled
continuation the wall is therefore about \SI{342}{\nano\meter} at
\SI{300}{\giga\pascal} and \SI{314}{\nano\meter} at
\SI{400}{\giga\pascal}. In a saturating continuation it remains near the
\SI{405}{\nano\meter} anchor value. These alternatives cannot be distinguished
by the present data. Moreover, pressure calibration itself becomes a material
systematic in this range: diamond-edge Raman scales have been demonstrated to
about \SI{370}{\giga\pascal}~\cite{Akahama2007} and extended to approximately
\SI{0.5}{\tera\pascal}, with substantial differences between competing scales
above \SI{300}{\giga\pascal}~\cite{Eremets2023}. A 300--400 GPa excitation
optimization is consequently feasible as an experiment, but not yet
identifiable as a unique calculation from the \SI{120}{\giga\pascal} record.

\subsection{Stress anisotropy and the status of the $\alpha$ correction}
\label{app:alpha}

The excitation wavelength can also depend on the stress tensor, not only on
its nominal pressure coordinate. For an NV axis along [111], we parameterize
an axisymmetric tensor by
\begin{equation}
  \boldsymbol{\sigma}(P,\alpha)
  = P\left[\alpha\,\mathbf{1}+(1-\alpha)\,\hat{\mathbf n}_{111}
      \hat{\mathbf n}_{111}^{\mathsf T}\right],
  \qquad \hat{\mathbf n}_{111}=\frac{(1,1,1)}{\sqrt{3}},
  \label{eq:alpha_tensor}
\end{equation}
where compression is positive. In the NV frame this is
$\operatorname{diag}(\alpha P,\alpha P,P)$, with mean stress
$\bar\sigma=(1+2\alpha)P/3$ and axial deviatoric excess
$\sigma_{\rm dev}= (1-\alpha)P$. Thus $\alpha=1$ is hydrostatic, whereas
smaller $\alpha$ represents increasing transverse stress deficit.

This distinction is important for interpreting the wavelength results. The
\SI{440.6}{\nano\meter} optimum reconstructed from Ref.~\cite{Ho2026} is a
hydrostatic optical-kernel result: $\alpha$ does not enter that kernel, so the
reported optimum is not an orientation-resolved prediction. The phenomenological
model used for the conditional stress analysis does include $\alpha$ in the
ZPL shift, calibrated to the ratio of the standard-culet and micropillar
responses reported by Hilberer \emph{et al.}~\cite{Hilberer2023}. At equal
nominal axial stress their reported values correspond approximately to
$\alpha=0.56$ for a standard flat culet and $\alpha=0.95$ for a micropillar;
the latter is close to quasihydrostatic loading. The earlier hydrostatic and
uniaxial-stress analysis of Doherty \emph{et al.}~\cite{Doherty2014} supplies
the physical basis for treating the ZPL as stress-sensitive.

The $\alpha$ correction is therefore only partial. It changes the ZPL
position, but the present data do not determine the corresponding
$\alpha$-dependence of the Huang--Rhys factor, phonon-sideband shape, linewidth,
charge-state yield, or the full absorption kernel. Consequently, the
stress-corrected wavelength is a conditional estimate, not a literature-
validated value of $\lambda_{\rm opt}(P,\alpha)$. In particular, substituting
$\bar\sigma$ for $P$ in the hydrostatic kernel would simply assume an
unmeasured mapping. The required validation is a wavelength-resolved
absorption/ODMR scan at several independently reconstructed $\alpha$ values,
with the local stress tensor measured simultaneously.

% =====================================================================
\section{Domain of validity of the reconstructed kernel}
\label{sec:kernel}

The worked examples of Secs.~\ref{sec:M} and~\ref{sec:X} descend from one
reconstruction of published curves---the exceptions are the reachability
estimate of Sec.~\ref{sec:M} and the drivers of Table~\ref{tab:drivers}, both
taken from published theory---and nothing in the framework's own logic can say
whether that reconstruction is sound. The check has to come from the
source.

\subsection{What the cross-figure check does and does not cover}
\label{sec:kernel:cross}

\begin{figure*}[t]
  \includegraphics[width=\textwidth]{figures/esa_figv_kernel_cross.pdf}
  \caption{The cross-figure check on the reconstructed kernel. Expt, the
  experimental gate; Ho, the theory gate of Ref.~\cite{Ho2026}. The absorption
  curve is reconstructed without reference to the independently digitized
  photoluminescence-yield curves, then compared with them at the two
  wavelengths those curves plot, (a) \SI{532}{\nano\meter} and
  (b) \SI{457}{\nano\meter}, giving fractional RMS residuals of $1.5\%$ and
  $0.4\%$ and reproducing the yield peaks at $17$ and \SI{88}{\giga\pascal}
  exactly; the vertical rule in each panel marks the published peak. The
  reconstruction carries one multiplicative scale and no shape freedom.
  (c) Those two residuals, and the two larger ones of $10.0\%$ and $10.2\%$
  obtained against the measured yields, each against its acceptance ceiling
  (red dash); the abscissa labels give the peak-position error and its ceiling
  in \si{\giga\pascal}. The figure validates an extraction of published
  curves; it is not an independent first-principles calculation.}
  \label{fig:kernel:cross}
\end{figure*}

The reconstruction reproduces the independently published pressure-dependent
photoluminescence yield with a pooled fractional RMS of $1.0\%$, correlation
$r = 1.00$, and peak pressures exact at $17$ and \SI{88}{\giga\pascal}, across
the full published \SIrange{0}{120}{\giga\pascal} axis
[Fig.~\ref{fig:kernel:cross}]. That is a strong check where it applies.

It probes exactly two photon energies, and both lie deep in the phonon
sideband throughout their audit windows: \SI{532}{\nano\meter}
(\SI{2.331}{\electronvolt}) over \SIrange{4.7}{51}{\giga\pascal}, sitting
$5.5$ to $2.2$ phonons above the ZPL, and \SI{457}{\nano\meter}
(\SI{2.713}{\electronvolt}) over \SIrange{51}{113.8}{\giga\pascal}, sitting
$8.1$ to $4.9$ phonons above it. \emph{The zero-phonon line is never sampled.}
The $1.0\%$ agreement must therefore never be cited as validating the ZPL
region, which is precisely where the reconstruction is known to be defective.

\subsection{Four internal checks}
\label{sec:kernel:internal}

\begin{figure*}[t]
  \includegraphics[width=\textwidth]{figures/esa_figv_kernel_internal.pdf}
  \caption{The four checks internal to the source, K1 to K4. ZPL, zero-phonon
  line; DWF, Debye-Waller factor. In (a) and (c) each curve is drawn in the
  color of the axis that scales it. (a) K1, the area normalization against
  pressure (left, black), falling to $0.914$ as the band ceases to return to
  zero inside the plotted window, which the weight remaining at the top of
  that window (right, red) tracks: $0.0\%$ of the peak at ambient pressure and
  $72.7\%$ at \SI{120}{\giga\pascal}. (b) K2, the apparent zero-phonon width in
  \si{\milli\electronvolt}, constant to $2.4\%$ over the full range and
  therefore resolution limited rather than physical. (c) K3, the kernel's own
  zero-phonon area against the published Debye-Waller factor (left, black),
  with the gray band marking $\pm5\%$ around parity and the green span the
  pressures below \SI{40}{\giga\pascal} over which the kernel ZPL may be used;
  and K4, the total coupling $S_{\mathrm{total}} = \Sabs + S_{JT}$, running
  from $3.83$ to $5.62$, whose Jahn-Teller-active remainder is $S_{JT}$
  (right, red).
  Panels (a) and (c) mark the two limits of use: the window edge, and the
  zero-phonon weight above \SI{40}{\giga\pascal}.}
  \label{fig:kernel:internal}
\end{figure*}

Four checks internal to the same source close that gap
[Fig.~\ref{fig:kernel:internal}].

\emph{K1: the blue edge of the analysis window is a figure limit.} The
published spectra are drawn to unit area, and the digitized curves confirm this
to $0.5\%$ at \SIrange{0}{60}{\giga\pascal}. Above that the integral falls, to
$0.914$ at \SI{120}{\giga\pascal}, because the band does not return to zero
inside the plotted window: absorption at \SIrange{3.0}{3.2}{\electronvolt} is
still $72.7\%$ of its own peak at \SI{120}{\giga\pascal}, against $0.0\%$ at
ambient pressure. The \SI{402}{\nano\meter} blue edge is therefore a property
of the figure, not of the defect, and the $\Delta N = -1$ step it generates in
Table~\ref{tab:ladder} is labeled an edge for that reason. Area-based
quantities are underestimated at high pressure by the missing tail.

\emph{K2: the ZPL width is resolution limited.} The apparent full width at half
maximum of the zero-phonon spike is \SIrange{2.53}{2.59}{\milli\electronvolt}
at all seven published pressures, constant to $2.4\%$ over
\SI{120}{\giga\pascal}. A physical ZPL broadens under compression---that is the
second term of Eq.~\eqref{eq:dDdP}. A width that does not move is the plotting
resolution. No ZPL width or broadening may be read from the reconstruction,
which is why Sec.~\ref{sec:X:bandwidth} replaces $\Gamma_\ZPL$ by the
excitation bandwidth rather than attempting to measure it.

\emph{K3: the kernel's ZPL carries a constant offset, not a trend.}
Integrating the background-subtracted spike gives the kernel's own
Debye-Waller factor. It is low against the published value by
\SIrange{9}{12}{\percent} at every one of the seven pressures, with no
monotone drift, and the two collapses over
\SIrange{0}{120}{\giga\pascal} agree to $1.3\%$---$\times6.08$ for the kernel
against $\times6.01$ for the published factor. The \emph{shape} of the
zero-phonon pressure dependence is therefore reproduced by the reconstruction;
what it misses is a scale. This is an independent validation in the region the
cross-figure check cannot reach, and it is why the zero-phonon critical value
of Table~\ref{tab:maxima} is quoted as a range rather than a number. Theorem~X
is nonetheless evaluated from the published Debye-Waller factor
(Table~\ref{tab:drivers}) rather than from the reconstruction, because there
the absolute scale is what enters.

An earlier raster extraction of the source's panel~(c) carried a constant
baseline offset and put this ratio at $0.96$ rising to $1.43$, which was read
as the reconstruction over-weighting the zero-phonon line above
\SI{40}{\giga\pascal}. That trend was an artifact of the extraction. The
vector reading used here reproduces the two endpoint values Ref.~\cite{Ho2026}
states in its own text, $2.2\%$ and $0.36\%$, to better than $1\%$.

\emph{K4: the two published panels are internally consistent.} The
Debye-Waller factor is not $\exp(-\Sabs)$, the ratio growing smoothly from
$2.11$ to $2.73$; this is expected, since $\Sabs$ is defined over the
Jahn-Teller-inactive modes only. The residual is the Jahn-Teller-active
coupling, $S_{\mathrm{JT}} = \ln[e^{-\Sabs}/\DWF]$, which runs from $0.745$ to
$1.005$ and gives $S_{\mathrm{total}}$ from $3.826$ to $5.619$. Two separately
digitized panels yielding a smooth, monotone derived quantity is itself
evidence that neither extraction is broken.

\subsection{Two regions that must not be used}
\label{sec:kernel:forbidden}

Two spectral regions of the reconstruction carry no data and must not be
quoted.

\emph{Below the zero-phonon line.} The lowest-energy real sample lies at
\SI{2.396}{\electronvolt} (\SI{517.4}{\nano\meter}), just below the
reconstructed ZPL at \SI{2.410}{\electronvolt} (\SI{514.5}{\nano\meter}); the
only other sample below that is a single baseline point at the figure's axis
edge. \SI{532}{\nano\meter} (\SI{2.331}{\electronvolt}) falls inside that
empty interval, so the ratio of \SI{457}{\nano\meter} to \SI{532}{\nano\meter}
sensitivity---nominally a factor $12$ in the optical limit---is produced by
interpolation, not by extracted data. Physically, absorption below the ZPL
vanishes at $T = 0$ and is thermally activated,
\begin{equation}
  \sabs(E_\gamma < E_\ZPL) \propto
  \exp\!\left[-\frac{E_\ZPL - E_\gamma}{k_BT}\right] ,
  \label{eq:hotband}
\end{equation}
with $E_\ZPL - E_{532} = \SI{79.5}{\milli\electronvolt}$ here, giving $4.6\times10^{-2}$ at \SI{300}{\kelvin} and $3.5\times10^{-5}$ at
\SI{90}{\kelvin}. The interpolated value $6.4\times10^{-3}$ lies between these
without deriving from either. Any comparison involving \SI{532}{\nano\meter} at
\SI{120}{\giga\pascal} is governed by hot-band absorption and temperature and
must be treated as temperature dependent, or the discussion restricted to
$\lambda \le \SI{517}{\nano\meter}$ where extracted data exist.

\emph{The zero-phonon peak heights.} Every ZPL spike in the source figure rises
to within one unit of the top of the axis---measured tops span only $0.92$
units across all seven pressures on an axis ending near $15$---so the spikes
are drawn as near-delta lines and the plot simply cuts them off. Their
reconstructed heights are a digitization artifact of a clipped feature. Taken
with K2, the only trustworthy zero-phonon information in the source figure is
its \emph{position}. This is why Sec.~\ref{sec:X} takes the ZPL weight from the
published Debye-Waller factor, and why $P^*$ is reported as a function of
bandwidth rather than as a number.

These restrictions reach Sec.~\ref{sec:M} unevenly. All four critical points of
Table~\ref{tab:maxima} lie at $\lambda \le \SI{514.5}{\nano\meter}$, inside the
region where extracted data exist, so their \emph{positions} are sound to the
node spacing---\SIrange{0.1}{0.5}{\nano\meter} at the zero-phonon point,
\SIrange{3}{10}{\nano\meter} in the sideband. One \emph{height} is affected,
and it is the one Theorem~M needs: $a(\lambda_\ZPL)$ is a zero-phonon peak
height, which is why Sec.~\ref{sec:M:example} carries it as a range and
propagates both branches of the resulting ladder. In an experiment the point is
moot, since T5 reads the $a_k$ from the measured spectrum.

\section{Reproducing the numbers}
\label{app:repro}

Every number in this paper is produced by the code in \texttt{code/} and pinned
by the regression suite.

\begin{itemize}
  \item \texttt{ho\_spectrum\_model.py}, \texttt{ho\_odmr\_sensitivity.py} ---
    the reconstructed kernel and the fixed-power sensitivity.
  \item \texttt{report\_120gpa\_sensitivity.py} --- the laser-line block of
    Table~\ref{tab:maxima}.
  \item \texttt{theory\_a1\_generalization.py} --- the coincidence condition,
    the exact degeneracy, the asymmetric $5\%$ band, and the T4 inference
    check of Sec.~\ref{sec:tests}. Its \texttt{v1\_envelope\_local\_maxima}
    counts the interior local maxima of the single-effective-mode envelope
    against pressure, on the grid, normalization and threshold that
    \texttt{Kernel.local\_maxima} applies to the reconstructed kernel; this is
    the one-against-four comparison of Sec.~\ref{sec:framework:kernel} and
    Sec.~\ref{sec:X:not}.
  \item \texttt{theory\_a2\_multiplicity.py} --- Theorem~M and Theorem~G:
    Tables~\ref{tab:maxima}, \ref{tab:ladder} and~\ref{tab:gauge}, the gauge
    degeneracy, and the invariance of the ladder. In
    \texttt{ladder\_is\_gauge\_invariant} each member of the gauge family is
    given its own half-scale $\Gamma$, so the five are probed at absolute
    intensities spanning a factor $16$ and the agreement of their ladders is a
    result rather than an identity.
  \item \texttt{dreau\_exponent.py} --- the splitting exponent $E = 3$,
    $\rho^* = 1/5$, and the reachability estimate of
    Eq.~\eqref{eq:reach}.
  \item \texttt{theory\_a3\_branch\_exchange.py},
    \texttt{figure\_validation.py} --- Theorem~X, the source-figure check, and
    the drivers of Table~\ref{tab:drivers} and the bandwidth-dependent
    $P^*(\Weff)$ curve of Fig.~\ref{fig:branch}(b).
  \item \texttt{kernel\_sanity\_checks.py} --- the checks K1--K4 of
    Appendix~\ref{sec:kernel:internal}.
  \item \texttt{repro\_yield.py} --- the cross-figure validation of
    Appendix~\ref{sec:kernel:cross}.
  \item \texttt{fig\_g\_gauge\_degeneracy.py},
    \texttt{fig\_m\_level\_set\_ladder.py},
    \texttt{fig\_x\_branch\_exchange.py},
    \texttt{esa\_figv\_kernel\_sanity.py} --- Figs.~\ref{fig:gauge},
    \ref{fig:levelset}, \ref{fig:branch}, \ref{fig:kernel:cross}
    and~\ref{fig:kernel:internal}. The last script draws the last two.
\end{itemize}

\begin{verbatim}
cd code
python report_120gpa_sensitivity.py
python theory_a1_generalization.py
python theory_a2_multiplicity.py
python theory_a3_branch_exchange.py
python figure_validation.py
python kernel_sanity_checks.py
python dreau_exponent.py
python fig_m_level_set_ladder.py
python fig_g_gauge_degeneracy.py
python fig_x_branch_exchange.py
python esa_figv_kernel_sanity.py
python -m pytest . -q
\end{verbatim}

The frozen record of the theory, including the errata that withdrew the
zero-phonon peak heights and the single-valued $P^*$, is
\path{docs/theory/theory_freeze_v3_ho_integrated.md}; the derivations are in
\path{docs/theory/theory_a1_numerical_execution.md},
\path{docs/theory/theory_a2_multiplicity.md} and
\path{docs/theory/theory_a3_branch_exchange.md}.

\bibliographystyle{apsrev4-2}
\bibliography{refs}

\end{document}
